\documentclass[11pt]{article}
\usepackage{tcee}
\addbibresource{3B41.bib}
\tcetema[Política comercial de la UE]{3.B.41}{La política comercial de la Unión Europea.}
\tcefecha{14/11/2026}
\begin{document}
\tceetitulo

\section*{Introducción}
\minutos{2,5}

\begin{itemize}
\item \textbf{\emph{Enganche}}:
  \begin{itemize}
  \item «Donde las dan, las toman», dice el refrán. El 2~de abril de 2025, Estados Unidos anunció un arancel «recíproco» del 20~\% para todo lo que le vendiera la Unión Europea. La Unión aprobó en julio contramedidas sobre 93.000~millones de euros de importaciones estadounidenses, pero no las aplicó: el 27~de julio pactó un techo del 15~\% para casi todas sus exportaciones a cambio de suprimir sus propios aranceles a los productos industriales de Estados Unidos\nocite{UE2025DeclaracionEEUU,CE2026Contramedidas}.
  \item Mientras tanto, en 2026 empezó a aplicar el acuerdo con el Mercosur, cerró las negociaciones con la India y con Australia y firmó el nuevo acuerdo con México\nocite{CE2026Informe}.
  \item La pregunta que recorre la exposición es: \emph{¿cómo defiende sus intereses una potencia comercial que vive de la apertura, cuando sus socios vuelven a las represalias?}
  \end{itemize}

\item \textbf{\emph{Relevancia}}:
  \begin{itemize}
  \item \emph{Es una potencia comercial}: sumando bienes y servicios, la UE es el mayor exportador del mundo, y es el segundo importador de bienes, y su red de acuerdos preferenciales es la mayor del mundo (44~acuerdos con 76~socios a finales de 2025, que cubren el 46,3~\% de su comercio exterior)\nocite{CE2026Informe}.
  \item \emph{Es la política más supranacional de la Unión}: es competencia exclusiva desde los Tratados de Roma, la Comisión negocia en nombre de los 27~y la UE habla con una sola voz en la Organización Mundial del Comercio (OMC).
  \item \emph{Para España}: desde 1986 su política comercial es la de la Unión \vertema{4.A.27}; el 61,8~\% de sus exportaciones de bienes va a otros Estados miembros, pero el 38,2~\% restante, unos 151.000~millones de euros en 2025, y el 43,9~\% de sus importaciones dependen de la política comercial común\nocite{Eurostat2026Espana} \vertema{4.A.22} \vertema{4.A.24}.
  \end{itemize}

\item \textbf{\emph{Contextualización}}:
  \begin{itemize}
  \item \emph{Qué es}: la \concepto{política comercial común} (PCC) es el conjunto de normas y decisiones con que la Unión regula los intercambios de bienes, servicios, propiedad intelectual e inversión directa con los países de fuera de la Unión (arts. 206~y 207~del Tratado de Funcionamiento de la Unión Europea, TFUE). Es la cara exterior de la unión aduanera: dentro, libre circulación \vertema{3.B.39}; frente a terceros, un arancel común y una sola política.
  \item \emph{Sus 2~vertientes}:
    \begin{itemize}
    \item la \concepto{política comercial autónoma}, que la Unión decide por sí sola: el arancel aduanero común, la defensa comercial, las preferencias unilaterales y los nuevos instrumentos de seguridad económica;
    \item la \concepto{política comercial convencional}, que pacta con otros: los acuerdos de la OMC y los acuerdos preferenciales.
    \end{itemize}
  \item \emph{Qué no trata esta exposición}: la base institucional de la Unión \vertema{3.B.36}, las reglas de la OMC \vertema{3.B.33} \vertema{3.B.34} y el marco político de las relaciones exteriores y de la cooperación al desarrollo \vertema{3.B.42}. Aquí, el instrumento comercial.
  \end{itemize}

\item \textbf{\emph{Problemática}}:
  \begin{itemize}
  \item ¿Por qué tiene la Unión una política comercial común y por qué la decide en exclusiva?
  \item ¿Con qué reglas se negocia, se aprueba y se ratifica un acuerdo comercial?
  \item ¿Qué instrumentos tiene para proteger a sus productores y para defenderse de la coerción?
  \item ¿Con quién comercia en condiciones preferentes, qué relación tiene con Estados Unidos y con China y qué retos le esperan?
  \end{itemize}

\item \textbf{\emph{Estructura}}:
  \begin{itemize}
  \item La exposición sigue el orden con que se estudia cualquier política de la Unión: por qué existe y quién la decide, cómo ha evolucionado, cómo se articula y qué retos tiene (figura~\ref{figu:mapa-tema}):
    \begin{itemize}
    \item \concepto{I. Las bases de la política comercial común} (\ref{sec:bases}): su relevancia; su justificación, con la teoría económica, la delimitación de competencias, el marco jurídico y la toma de decisiones; y su evolución, de Roma a la seguridad económica.
    \item \concepto{II. La política comercial autónoma} (\ref{sec:autonoma}): la primera mitad de su articulación, lo que la Unión decide sola: el arancel y el TARIC, la defensa comercial, las preferencias generalizadas y los nuevos instrumentos.
    \item \concepto{III. La política comercial convencional} (\ref{sec:convencional}): la segunda mitad, lo que pacta: la OMC, la red de acuerdos preferenciales y las relaciones con Estados Unidos y con China; y, para cerrar, sus retos de futuro.
    \end{itemize}
  \end{itemize}
\end{itemize}

\begin{esquema}[unbreakable]
\textbf{3.B.41. La política comercial de la Unión Europea}\\[2pt]
Introducción\esqtiempo{2,5}\\
I. Las bases de la política comercial común\esqtiempo{7}\\
\hspace*{1.5em}1. Relevancia\\
\hspace*{1.5em}2. Aproximación de teoría económica\\
\hspace*{1.5em}3. Delimitación de competencias\\
\hspace*{1.5em}4. Marco jurídico y toma de decisiones\\
\hspace*{1.5em}5. La evolución: de Roma a la seguridad económica\\
II. La política comercial autónoma\esqtiempo{9,5}\\
\hspace*{1.5em}1. La política arancelaria\\
\hspace*{1.5em}2. Las medidas de defensa comercial\\
\hspace*{1.5em}3. El Sistema de Preferencias Generalizadas\\
\hspace*{1.5em}4. Los instrumentos de seguridad económica y sostenibilidad\\
III. La política comercial convencional\esqtiempo{10}\\
\hspace*{1.5em}1. Los acuerdos multilaterales y plurilaterales de la OMC\\
\hspace*{1.5em}2. Los acuerdos comerciales preferenciales\\
\hspace*{1.5em}3. Acuerdos en negociación y pendientes de ratificación\\
\hspace*{1.5em}4. Otras relaciones comerciales: Estados Unidos, China y Rusia\\
\hspace*{1.5em}5. Retos de futuro\\
Conclusión\esqtiempo{1}
\end{esquema}

\begin{anotaciones}
\textbf{Mapa de la exposición.} Cada bloque responde a una pregunta y deja abierta la siguiente; a la izquierda, el elemento de la estructura común de las políticas de la Unión que cubre.
\begin{infografiacaja}
\caption{Mapa de la exposición: la pregunta de cada bloque y el hilo que los une}\label{figu:mapa-tema}
\grafico[0.86\linewidth]{mapa-tema}
\fuente{elaboración propia.}
\end{infografiacaja}
\end{anotaciones}

\begin{anotaciones}
\textbf{Siglas.} AAC: arancel aduanero común; ALC: acuerdo de libre comercio; NMF: nación más favorecida; SA: Sistema Armonizado; NC: Nomenclatura Combinada; TARIC: arancel integrado de la Unión. La notación de los 3~modelos está en sus anexos.
\end{anotaciones}


% ==========================================================================
\section{Las bases de la política comercial común}\label{sec:bases}
\minutos{7}

\begin{itemize}
\item Empezamos por los fundamentos: cuánto pesa la Unión en el comercio mundial, por qué tiene sentido una política comercial única y con qué reglas se hace.
\item \emph{Conclusión anticipada}: la PCC es la consecuencia lógica de la unión aduanera y el instrumento que convierte a 27~economías medianas en una potencia negociadora; su eficacia depende de cómo se reparte el poder entre la Comisión, el Consejo, el Parlamento y los Estados.
\end{itemize}

% --------------------------------------------------------------------------
\subsection{Relevancia}\label{sec:relevancia}

\begin{itemize}
\item \emph{El comercio forma parte del ADN de la Unión}: la integración europea empezó con el comercio (la unión aduanera de 1968) y sigue siendo una de las economías más abiertas entre las grandes: en 2025 exportó 2,65~billones de euros en bienes y 1,58~en servicios fuera de la Unión, y compró 2,54~y 1,43~billones\nocite{Eurostat2026Socios,Eurostat2026Servicios}.
\item \emph{Sus socios} (gráfico~\ref{fig:socios}):
  \begin{itemize}
  \item \emph{Estados Unidos} es el primer socio: le vende el 21~\% de sus exportaciones de bienes, con un superávit de casi 200.000~millones de euros, que se compensa en buena parte con el déficit en servicios (178.000~millones)\nocite{Eurostat2026Servicios};
  \item \emph{China} es el segundo socio y el primer proveedor: la UE le compra 574.000~millones y le vende 199.000, con un déficit de 375.000~millones, el mayor con cualquier país\nocite{Eurostat2026Socios};
  \item siguen el \emph{Reino Unido}, \emph{Suiza} y \emph{Turquía}, los 3~con acuerdos preferenciales muy profundos.
  \end{itemize}
\item \emph{Sus intereses}: es ofensiva en manufacturas, servicios, contratación pública, indicaciones geográficas y reglas, y defensiva en la agricultura y en los sectores con exceso de capacidad mundial (acero, química).
\end{itemize}

\begin{figure}[H]
\centering
\caption{Comercio de bienes de la UE con sus 16~primeros socios, 2025 (miles de millones de euros)}\label{fig:socios}
\grafico[0.92\textwidth]{comercio-socios}
\fuente{Eurostat, \emph{Extra-EU trade by partner} (ext\_lt\_maineu), datos de 2025\nocite{Eurostat2026Socios}; situación de los acuerdos en noviembre de 2026, Comisión Europea\nocite{CE2026Negociaciones}.}
\end{figure}

\begin{anotaciones}
\textbf{El comercio de servicios por socios, 2025} (miles de millones de euros, exportaciones/importaciones de la UE)\nocite{Eurostat2026Servicios}: Estados Unidos, 343/522; Reino Unido, 309/249; Suiza, 172/102; China, 67/46; Japón, 45/24; la India, 29/38; Canadá, 29/20; Brasil, 24/10. La UE tiene superávit en servicios con casi todos, salvo con Estados Unidos (servicios digitales, propiedad intelectual y financieros) y con la India (informática). Sumando bienes y servicios, el superávit con Estados Unidos es de unos 18.000~millones de euros, mucho menor que el que se ve solo en bienes: es uno de los argumentos que la UE usó en la negociación de 2025.
\end{anotaciones}

% --------------------------------------------------------------------------
\subsection{Aproximación de teoría económica}\label{sec:teoria}

\subsubsection{¿Por qué una política comercial?}\label{sec:porque}

\begin{itemize}
\item \emph{El punto de partida es el librecambio}: los modelos de comercio dicen que abrirse da ganancias, por la ventaja comparativa (\autor{Ricardo}, \autor{Heckscher}-\autor{Ohlin}), por las economías de escala y la variedad (nueva teoría del comercio) y por la selección de las empresas más productivas (\autor{Melitz}, 2003) \vertema{3.B.5} \vertema{3.B.6}; y, a largo plazo, por la difusión de ideas \vertema{3.B.9}.
\item \emph{Pero hay excepciones}, que dan razón de ser a una política comercial activa \vertema{3.B.7} \vertema{3.B.8}:
  \begin{itemize}
  \item el \emph{arancel óptimo}: un país grande, como la UE, puede mejorar su relación real de intercambio (RRI) con un arancel, si los demás no responden;
  \item la \emph{industria naciente} y la \emph{política comercial estratégica}, que trasladan rentas a las empresas propias en mercados de oligopolio;
  \item la \emph{protección frente a la competencia desleal} (dumping y subvenciones) y frente a las oleadas de importaciones;
  \item y, cada vez más, la \emph{seguridad}: depender de un solo proveedor de bienes críticos es un riesgo que el mercado no valora.
  \end{itemize}
\item \emph{La integración regional es un segundo óptimo}: si el librecambio mundial no es posible, una unión aduanera puede crear comercio (sustituye producción interna cara por la de un socio eficiente) o desviarlo (sustituye a un proveedor de fuera eficiente por un socio menos eficiente), según el análisis de \textcite{Viner1950Customs} \vertema{3.B.10}. Cuantos más socios y más bajo el arancel exterior común, más creación y menos desviación: una buena razón para que la UE combine su unión aduanera con un arancel exterior bajo y una red de acuerdos.
\end{itemize}

\subsubsection{¿Por qué una competencia exclusiva?}\label{sec:exclusiva}

\begin{itemize}
\item \emph{Es una consecuencia de la unión aduanera}: con un arancel común, una política comercial nacional es imposible; cada mercancía que entra por un puerto circula libremente por toda la Unión. Si cada Estado pactara sus propias rebajas, el comercio entraría por el más abierto (\concepto{desviación del tráfico}, \emph{trade deflection}) y haría falta volver a las aduanas interiores.
\item \emph{Reduce los costes de transacción}: un solo arancel, unas solas reglas de origen y una sola negociación en lugar de 27.
\item \emph{Aumenta el poder de negociación}: un mercado de 450~millones de consumidores pesa más que cualquiera de sus Estados por separado, en la OMC y en los acuerdos bilaterales. Es la intuición de «la unión hace la fuerza», que el modelo siguiente matiza.
\item \emph{Encaja con la teoría del federalismo}: la política comercial tiene grandes externalidades entre Estados y preferencias relativamente parecidas, que es cuando conviene centralizar una competencia \vertema{3.B.36} \vertema{4.B.14}.
\end{itemize}

\subsubsection{La economía política de la «sola voz»}\label{sec:meunier}

\begin{itemize}
\item \emph{Hablar con una sola voz no basta para negociar mejor}: \textcite{Meunier2005Trading} muestra que el resultado de una negociación de la UE depende de 3~cosas:
  \begin{itemize}
  \item la \emph{regla de voto} en el Consejo: unanimidad o mayoría cualificada;
  \item la \emph{delegación} en la Comisión: un mandato estrecho y vigilado o uno amplio;
  \item y el \emph{contexto}: si es el socio el que quiere abrir el mercado europeo (caso \emph{conservador}, como la agricultura) o la UE la que quiere abrir el del socio (caso \emph{reformista}, como la contratación pública).
  \end{itemize}
\item \emph{La idea}: el socio solo consigue lo que acepta el Estado cuyo voto decide (el \concepto{Estado pivote}). Con unanimidad decide el más reacio, y la UE es un negociador duro en lo defensivo; con mayoría cualificada decide un Estado intermedio, y la UE puede cerrar acuerdos ofensivos. Es la aplicación a la Unión de las «manos atadas» de \autor{Thomas Schelling}\footnote{\autor{Thomas Schelling} recibió el Premio Nobel de Economía en 2005, con \autor{Robert Aumann}, «por haber mejorado nuestra comprensión del conflicto y la cooperación mediante el análisis de la teoría de juegos»\nocite{Nobel2005AumannSchelling}.} (1960)\nocite{Schelling1960Strategy} y de los juegos a 2~niveles de \textcite{Putnam1988Diplomacy}.
\item \emph{Lo que explica}: por qué la UE es dura en lo defensivo y lenta en lo ofensivo, por qué el socio divide a los Estados cuando no hay competencia común (los «cielos abiertos» de la aviación en los años 90) y por qué el acuerdo con el Mercosur se partió en 2~para aprobar la parte comercial por mayoría cualificada. El modelo completo, con su gráfico, está en el anexo~\ref{anexo:meunier}.
\end{itemize}

% --------------------------------------------------------------------------
\subsection{Delimitación de competencias}\label{sec:delimitacion}

\begin{itemize}
\item \emph{Competencia exclusiva}: la política comercial común es una de las 5~competencias exclusivas de la Unión (art. 3.1.e TFUE), junto con la unión aduanera: solo la Unión puede legislar y celebrar acuerdos; los Estados, solo si la Unión los habilita (art. 2.1). También es exclusiva la celebración de acuerdos internacionales cuando puede afectar a normas comunes (art. 3.2). El reparto completo de competencias de los arts. 2~a 6~del TFUE está en el anexo~\ref{anexo:competencias} \vertema{3.B.36}.
\item \emph{Su ámbito}: el comercio de mercancías y servicios, los aspectos comerciales de la propiedad intelectual y la \concepto{inversión extranjera directa} (IED) con los países de fuera de la Unión (art. 207.1).
\item \emph{Política autónoma frente a convencional}: la primera es lo que la Unión decide sola (bloque~\ref{sec:autonoma}); la segunda, lo que pacta con otros (bloque~\ref{sec:convencional}).
\item \emph{Lo que no es PCC}:
  \begin{itemize}
  \item la \emph{inversión de cartera} y la solución de diferencias entre inversores y Estados son competencias compartidas (dictamen 2/15, \ref{sec:dictamen});
  \item la \emph{promoción comercial} y el apoyo financiero a la exportación siguen siendo nacionales (el ICEX, la CESCE o el FIEM en España) \vertema{4.A.24} \vertema{4.A.25}, aunque la UE vela por que no falseen la competencia y aplica el Consenso de la OCDE sobre créditos a la exportación con apoyo oficial \vertema{3.B.8};
  \item la \emph{agricultura} tiene su propia política, aunque sus aranceles y contingentes forman parte del AAC \vertema{3.B.38};
  \item las \emph{sanciones} (medidas restrictivas) son de la política exterior y de seguridad común y se deciden por unanimidad (art. 215~TFUE), aunque usan instrumentos comerciales.
  \end{itemize}
\end{itemize}

% --------------------------------------------------------------------------
\subsection{Marco jurídico y toma de decisiones}\label{sec:juridico}

\subsubsection{Base jurídica y objetivos}\label{sec:base-juridica}

\begin{itemize}
\item \emph{Base jurídica}: los arts. 206~y 207~del TFUE (la PCC) y el art. 218 (el procedimiento de los acuerdos internacionales).
\item \emph{Objetivos} (art. 206~TFUE): contribuir al desarrollo armonioso del comercio mundial, a la supresión progresiva de las restricciones a los intercambios y a la IED y a la reducción de las barreras arancelarias y de otro tipo. Se ejerce en el marco de los principios y objetivos de la acción exterior (art. 21~TUE): derechos humanos, desarrollo sostenible, multilateralismo.
\end{itemize}

\subsubsection{El dictamen 2/15 y la nueva arquitectura de los acuerdos}\label{sec:dictamen}

\begin{itemize}
\item \emph{La pregunta}: tras Lisboa, ¿podía la Unión celebrar sola un acuerdo «de nueva generación» como el de Singapur, que trata aranceles, servicios, contratación pública, propiedad intelectual, inversión y desarrollo sostenible?
\item \emph{El dictamen 2/15} (Tribunal de Justicia, 16~de mayo de 2017)\nocite{TJUE2017Dictamen215}:
  \begin{itemize}
  \item casi todo el acuerdo es competencia exclusiva de la Unión: el comercio de mercancías y servicios (incluido el transporte), la contratación pública, la propiedad intelectual, la competencia, la IED y también el capítulo de desarrollo sostenible, porque está ligado al comercio;
  \item pero 2~materias son compartidas con los Estados: la \emph{inversión de cartera} (no directa) y la \emph{solución de diferencias entre inversores y Estados}, porque saca los litigios de los tribunales nacionales.
  \end{itemize}
\item \emph{La respuesta, la nueva arquitectura} (Conclusiones del Consejo del 22~de mayo de 2018)\nocite{Consejo2018Arquitectura}:
  \begin{itemize}
  \item como regla, el acuerdo de libre comercio va \concepto{solo UE} (\emph{EU-only}): lo celebran el Consejo, por mayoría cualificada, y el Parlamento Europeo, sin ratificaciones nacionales;
  \item la protección de inversiones va en un acuerdo aparte, \concepto{mixto}, que ratifican los 27;
  \item se ha aplicado a Japón, Singapur, Vietnam, Nueva Zelanda, Chile, México y el Mercosur, en estos 3~últimos con un \emph{acuerdo comercial interino} que se aplica mientras se ratifica el acuerdo de asociación completo.
  \end{itemize}
\item \emph{Por qué importa}: un acuerdo mixto puede quedar años a medias. El CETA con Canadá se aplica de forma provisional desde 2017 y en noviembre de 2026 todavía no lo habían ratificado 9~de los 27~Estados (Bélgica, Bulgaria, Chipre, Francia, Grecia, Hungría, Irlanda, Italia y Polonia; en 2024 lo rechazó el Senado francés), por lo que su capítulo de inversiones no se aplica\nocite{CE2026CETA}.
\end{itemize}

\subsubsection{La toma de decisiones: cómo se negocia y se celebra un acuerdo}\label{sec:procedimiento}

\begin{itemize}
\item \emph{Las medidas autónomas} (el arancel, la defensa comercial, el SPG) se adoptan por reglamentos del Parlamento Europeo y del Consejo, por el procedimiento legislativo ordinario (art. 207.2), y la Comisión las ejecuta con actos de ejecución, bajo el control de comités de los Estados (procedimiento de examen).
\item \emph{Los acuerdos} siguen los arts. 207.3~y 218~TFUE (figura~\ref{figu:procedimiento}):
  \begin{itemize}
  \item la \emph{Comisión} recomienda abrir negociaciones; el \emph{Consejo} la autoriza y le da directrices;
  \item la Comisión negocia en consulta con el \concepto{Comité de Política Comercial} (el comité especial del art. 207.3, de representantes de los Estados) e informa por igual al Parlamento;
  \item el Consejo decide la firma; el \emph{Parlamento Europeo} da su aprobación, sin poder enmendar; el Consejo celebra el acuerdo;
  \item si el acuerdo es mixto, además, lo ratifican los 27~según sus normas constitucionales.
  \end{itemize}
\item \emph{Mayorías en el Consejo} (art. 207.4): cualificada, salvo unanimidad para los servicios, la propiedad intelectual y la IED cuando el acuerdo incluya disposiciones que exijan unanimidad para las normas internas, y para los servicios culturales y audiovisuales (si afectan a la diversidad cultural) y los sociales, educativos y sanitarios (si pueden perturbar su organización nacional).
\end{itemize}

\begin{infografia}[H]
\centering
\caption{Cómo se negocia y se celebra un acuerdo comercial de la UE (arts. 207~y 218~TFUE)}\label{figu:procedimiento}
\grafico[0.95\textwidth]{procedimiento-acuerdos}
\fuente{elaboración propia a partir de los arts. 207~y 218~del TFUE\nocite{UE2016TFUE} y de las Conclusiones del Consejo de 2018\nocite{Consejo2018Arquitectura}.}
\end{infografia}

\subsubsection{Las reglas de la OMC como marco}\label{sec:omc-marco}

\begin{itemize}
\item La UE es miembro de la OMC desde su creación (1995), junto con sus 27~Estados; la Comisión habla en nombre de todos \vertema{3.B.33}. Su política comercial tiene que respetar 2~excepciones a la cláusula de \concepto{nación más favorecida} (NMF):
  \begin{itemize}
  \item \emph{art. XXIV del GATT}: las uniones aduaneras y zonas de libre comercio se permiten si liberalizan «lo esencial de los intercambios» (en la práctica, en torno al 90~\%) y no elevan la protección frente a terceros. Si una ampliación sube un arancel consolidado, la UE tiene que compensar: cuando Croacia adoptó el arancel común en 2013, China pidió y obtuvo un contingente de ajos;
  \item la \emph{cláusula de habilitación} (1979): permite preferencias no recíprocas a los países en desarrollo, la base del SPG (\ref{sec:spg}).
  \end{itemize}
\end{itemize}

% --------------------------------------------------------------------------
\subsection{La evolución: de Roma a la seguridad económica}\label{sec:estrategia}

\subsubsection{El alcance: de Roma a Lisboa}\label{sec:alcance}

\begin{itemize}
\item \emph{El ámbito de la PCC ha crecido con cada Tratado} (figura~\ref{figu:cronologia}):
  \begin{itemize}
  \item \emph{Roma} (1957): una política comercial común basada en «principios uniformes», sobre todo de las mercancías. El Tribunal de Justicia confirmó en 1975 (dictamen 1/75) que la competencia era exclusiva.
  \item \emph{Dictamen 1/94}: el Tribunal resolvió que la Comunidad tenía competencia exclusiva para los acuerdos de mercancías de la Ronda Uruguay, pero compartida para la mayor parte de los servicios y de la propiedad intelectual: los acuerdos de la OMC fueron \emph{mixtos}.
  \item \emph{Ámsterdam} (1997) y \emph{Niza} (2001): extendieron la PCC a los servicios y a los aspectos comerciales de la propiedad intelectual, con excepciones y unanimidad.
  \item \emph{Lisboa} (2009): la PCC cubre todo el comercio de servicios, los aspectos comerciales de la propiedad intelectual y la IED; las medidas autónomas se adoptan por el procedimiento legislativo ordinario y el Parlamento Europeo tiene que aprobar todos los acuerdos.
  \end{itemize}
\end{itemize}

\begin{infografia}[H]
\centering
\caption{La política comercial común, 1957-2026}\label{figu:cronologia}
\grafico[0.98\textwidth]{cronologia}
\fuente{elaboración propia a partir de los Tratados\nocite{UE2016TFUE} y de la Comisión Europea\nocite{CE2021TradePolicyReview,CE2023SeguridadEconomica}.}
\end{infografia}

\subsubsection{La estrategia: del multilateralismo a la seguridad económica}\label{sec:estrategia-etapas}
\begin{itemize}
\item \emph{Multilateralismo primero} (1957-2006): la Comunidad negoció en las rondas del GATT y defendió la OMC, ofensiva en manufacturas y en los «temas de Singapur» y defensiva en la agricultura. Hasta los 2000 solo tenía acuerdos preferenciales con sus vecinos y con las antiguas colonias.
\item \emph{Europa Global} (\emph{Global Europe}, 2006): ante el bloqueo de Doha, la UE se lanza a los acuerdos bilaterales con economías grandes y dinámicas (Corea del Sur, la India, la ASEAN), con criterios de tamaño del mercado y de barreras.
\item \emph{Comercio para todos} (\emph{Trade for All}, 2015)\nocite{CE2015TradeForAll}: tras las protestas contra el TTIP, la política comercial tiene que ser eficaz, transparente y basada en valores (desarrollo sostenible, derechos humanos); los textos negociadores se publican.
\item \emph{Revisión de 2021: abierta, sostenible y firme}\nocite{CE2021TradePolicyReview}: introduce la \concepto{autonomía estratégica abierta}: seguir abierta, pero con capacidad de decidir por sí misma y de defenderse. Sus prioridades: reformar la OMC, apoyar la transición verde y digital, diversificar las cadenas de suministro y ser más firme ante las prácticas desleales.
\item \emph{La seguridad económica} (2023-2026):
  \begin{itemize}
  \item la Estrategia Europea de Seguridad Económica (junio de 2023)\nocite{CE2023SeguridadEconomica} la ordena en 3~ejes: \emph{promover} la competitividad propia, \emph{proteger} frente a los riesgos (cadenas de suministro, tecnología, coerción, seguridad física) y \emph{asociarse} con socios afines;
  \item la Comunicación de diciembre de 2025\nocite{CE2025SeguridadEconomica} la convierte en doctrina: usar todos los instrumentos, sea cual sea su fin original, y coordinar a los Estados;
  \item en octubre de 2026, el Consejo Europeo pidió a la Comisión «reforzar urgentemente» la caja de herramientas frente a los desequilibrios mundiales y al exceso de capacidad\nocite{CEU2026ConclusionesOctubre}.
  \end{itemize}
\item \emph{Recapitulación del bloque~\ref{sec:bases}}: la PCC nace de la unión aduanera, se justifica por el poder de negociación y su alcance lo ha fijado el Tribunal de Justicia. ¿Con qué instrumentos actúa la Unión cuando decide sola?
\end{itemize}

% ==========================================================================
\section{La política comercial autónoma}\label{sec:autonoma}
\minutos{9,5}

\begin{itemize}
\item Una vez vistos los fundamentos, pasamos a lo que la Unión decide sola: el arancel, la defensa comercial, las preferencias unilaterales y la nueva caja de herramientas de la seguridad económica.
\item \emph{Conclusión anticipada}: el arancel ya protege poco, salvo en la agricultura; la protección se ha desplazado a la defensa comercial y, desde 2020, a instrumentos nuevos que ya no miran solo al precio, sino a la coerción, a la seguridad y a la sostenibilidad.
\end{itemize}

% --------------------------------------------------------------------------
\subsection{La política arancelaria}\label{sec:arancel}

\subsubsection{El arancel aduanero común}\label{sec:aac}

\begin{itemize}
\item \emph{Origen}: la unión aduanera se completó el 1~de julio de 1968, 18~meses antes de lo previsto, con un \concepto{arancel aduanero común} (AAC) que se fijó como la media aritmética de los aranceles de las 4~zonas aduaneras de los Seis (Francia, Alemania, Italia y el Benelux). Los Estados que se adhieren lo adoptan tras un periodo transitorio \vertema{3.B.39}.\footnote{Forman la unión aduanera los 27~y, por acuerdos con la UE, Mónaco y, para la mayoría de los productos, Andorra, San Marino y Turquía (\ref{sec:preferenciales}). Algunos territorios de Estados miembros están fuera: en España, Ceuta y Melilla; las Canarias están dentro de la unión aduanera, pero fuera del territorio del IVA.}
\item \emph{Cómo es}:
  \begin{itemize}
  \item se aplica a cada mercancía según su código de la \emph{Nomenclatura Combinada} (8~dígitos), y el \emph{TARIC} (10~dígitos) le suma las demás medidas (\ref{sec:taric});
  \item es bajo: el arancel NMF aplicado medio es del 5,0~\% en media simple (2024) y del 2,8~\% ponderado por las importaciones; pero es muy desigual, del 10,5~\% en los productos agrícolas frente al 4,1~\% en los demás (gráfico~\ref{fig:aranceles})\nocite{OMC2025PerfilesArancelarios};
  \item quedan \concepto{crestas arancelarias} en los lácteos, el azúcar, la carne, las bebidas, el pescado, la ropa y los automóviles (10~\%), y los productos agrícolas tienen muchos aranceles específicos y contingentes.
  \end{itemize}
\item \emph{Su recaudación}, los \emph{recursos propios tradicionales} del presupuesto de la Unión: las aduanas nacionales la recaudan y se quedan el 25~\% como gastos de recaudación \vertema{3.B.37}.\footnote{De ahí el «efecto Róterdam»: los derechos los ingresa el Estado por donde entra la mercancía (los Países Bajos o Bélgica), aunque el producto se consuma en otro, y las estadísticas de comercio de esos Estados se inflan con mercancía en tránsito.}
\end{itemize}

\begin{figure}[H]
\centering
\caption{Arancel NMF aplicado medio de la UE por grupos de productos, 2024 (\%)}\label{fig:aranceles}
\grafico[0.8\textwidth]{aranceles-ue}
\fuente{OMC, ITC y UNCTAD, \emph{World Tariff Profiles 2025}, ficha de la Unión Europea\nocite{OMC2025PerfilesArancelarios}. Media simple de los tipos NMF aplicados en cada grupo; no incluye las preferencias de los acuerdos ni los derechos de defensa comercial.}
\end{figure}

\begin{itemize}
\item \emph{Media simple o media ponderada}: las 2~medias dicen cosas distintas, y conviene saber cuál se cita.
  \begin{itemize}
  \item La \concepto{media simple} da el mismo peso a cada línea arancelaria: describe la estructura del arancel, con sus crestas, sea cual sea lo que se importa.
  \item La \concepto{media ponderada} pesa cada línea por lo que se importa: dice cuánto se paga de media por lo que de verdad entra.
  \item La ponderada \emph{infravalora la protección}: los productos con aranceles altos se importan poco precisamente porque el arancel es alto, y pesan poco en la media; un arancel prohibitivo, que no deja entrar nada, pesa cero. Por eso la ponderada de la UE (2,8~\%) es mucho menor que la simple (5,0~\%), sobre todo en la agricultura.
  \end{itemize}
\item \emph{Su evolución} (gráfico~\ref{fig:aranceles-evolucion}): desde 1988 la media simple del arancel NMF ha bajado de algo más del 6~\% a algo más del 5~\%, y la ponderada, del 4,7~\% al 2,7~\%\nocite{BM2026Aranceles}.
  \begin{itemize}
  \item La subida de 1995 es la \emph{arancelización} de la Ronda Uruguay: las barreras no arancelarias agrícolas (cuotas y exacciones variables) se convirtieron en aranceles, a menudo muy altos, que después se fueron rebajando.
  \item La distancia entre las 2~medias ha crecido: la protección que queda se concentra en productos que se importan poco.
  \item Con las preferencias de los acuerdos y del SPG, el arancel que de verdad se aplica es aún menor: un 1,3~\% ponderado en 2022\nocite{BM2026Aranceles}. El arancel ya no es la principal barrera: lo son las normas, la defensa comercial y los nuevos instrumentos.
  \end{itemize}
\end{itemize}

\begin{figure}[H]
\centering
\caption{Evolución del arancel medio de la UE, 1988-2022 (\%): media simple y ponderada}\label{fig:aranceles-evolucion}
\grafico[0.95\textwidth]{aranceles-evolucion}
\fuente{Banco Mundial, \emph{World Development Indicators}, con datos de WITS-TRAINS, todos los productos\nocite{BM2026Aranceles}. NMF: el arancel de la nación más favorecida; aplicado: el más bajo que se aplica a cada importación, con las preferencias. La media ponderada usa las importaciones de cada producto. Los saltos de un solo año conviene leerlos con cautela: dependen de la cobertura de las líneas y del cálculo de los equivalentes ad valorem de los derechos específicos. El último dato de la OMC (2024) es del 5,0~\% simple y del 2,8~\% ponderado\nocite{OMC2025PerfilesArancelarios}.}
\end{figure}

\subsubsection{La nomenclatura: del Sistema Armonizado al TARIC}\label{sec:taric}

\begin{itemize}
\item \emph{Por qué importa}: el arancel no se aplica a «las bicicletas» en general, sino a un código; clasificar bien una mercancía decide cuánto paga y qué medidas la alcanzan.
\item \emph{La estructura del código} (tabla~\ref{tab:taric}): cada nivel añade 2~dígitos al anterior.
  \begin{itemize}
  \item \emph{2 dígitos}: el \emph{capítulo} del \concepto{Sistema Armonizado} (SA) de la Organización Mundial de Aduanas, la nomenclatura común a más de 200~países (97~capítulos);
  \item \emph{4 dígitos}: la \emph{partida} del SA;
  \item \emph{6 dígitos}: la \emph{subpartida} del SA, el último nivel común en todo el mundo;
  \item \emph{8 dígitos}: la subpartida de la \concepto{Nomenclatura Combinada} (NC), la de la UE, que fija el derecho de terceros países (el NMF) y sirve para las estadísticas de comercio; se publica cada año en el anexo I del Reglamento 2658/87;
  \item \emph{10 dígitos}: la subpartida \concepto{TARIC} (arancel integrado), que no cambia el derecho de terceros países pero permite aplicar las demás medidas de la Unión: preferencias, contingentes, suspensiones, derechos antidumping y compensatorios, prohibiciones y controles\nocite{CE2026TARIC};
  \item \emph{códigos adicionales} de 4~caracteres: distinguen, dentro de un mismo código, por ejemplo al exportador al que se aplica cada derecho antidumping.
  \end{itemize}
\end{itemize}

\begin{table}[H]
\centering
\caption{Un ejemplo de clasificación: la bicicleta con cojinetes de bolas en el TARIC}\label{tab:taric}
\footnotesize
\begin{tabular}{>{\raggedright\arraybackslash}m{2.6cm}>{\raggedright\arraybackslash}m{2.4cm}>{\raggedright\arraybackslash}m{9.6cm}}
\toprule
\rowcolor{tceeazulpalido} \textbf{Nivel} & \textbf{Código} & \textbf{Descripción y medida}\\
\midrule
Capítulo (SA) & 87 & Vehículos automóviles, tractores, velocípedos y demás vehículos terrestres, sus partes y accesorios\\
Partida (SA) & 8712 & Bicicletas y demás velocípedos (incluidos los triciclos de reparto), sin motor\\
Subpartida (SA) & 8712 00 & La misma (la partida no se divide en el SA)\\
Nomenclatura Combinada & 8712 00 30 & Bicicletas con cojinetes de bolas: \textbf{derecho de terceros países del 14~\%}\\
TARIC & 8712 00 30 10 & Procedentes de Indonesia, Malasia, Sri Lanka o Túnez (antidumping extendido por elusión)\\
 & 8712 00 30 20 & Procedentes de Camboya, Pakistán o Filipinas (ídem)\\
 & 8712 00 30 90 & Las demás\\
Código adicional & B999 & «Las demás» empresas de China: \textbf{antidumping del 48,5~\%}; otros códigos identifican a los exportadores con trato propio\\
\bottomrule
\end{tabular}
\fuente{Comisión Europea, consulta TARIC del código 8712~00~30~con fecha de referencia 19/11/2026\nocite{CE2026TARICBicicletas}.}
\end{table}

\begin{itemize}
\item \emph{Cuánto paga una bicicleta según su origen} (consulta del 19/11/2026)\nocite{CE2026TARICBicicletas}:
  \begin{itemize}
  \item de \emph{China}: el 14~\% del derecho de terceros países más un \emph{derecho antidumping definitivo del 48,5~\%} (Reglamento de Ejecución 2025/2146, en vigor desde el 24~de octubre de 2025, que prorroga una medida que existe desde 1993)\nocite{CE2026TARICBicicletas};
  \item de \emph{Indonesia}: el SPG le rebaja el derecho al 10,5~\%, pero el antidumping chino se le extiende por elusión (48,5~\%), salvo a los fabricantes que demostraron no reexpedir bicicletas chinas, que tienen su propio código adicional: es para eso para lo que existe la subpartida TARIC «10»;
  \item de \emph{Vietnam}: el 0~\%, por su acuerdo de libre comercio con la UE, que liberalizó las bicicletas el 1~de enero de 2025\nocite{CE2026TARICBicicletas};
  \item y, si llega en un envío de escaso valor de comercio electrónico, además el derecho a tanto alzado de 3~euros y la tasa de gestión de 2~euros (\ref{sec:modulacion}).
  \end{itemize}
\item Un mismo producto puede pagar, por tanto, del 0~\% al 62,5~\% según su origen: el arancel común es el punto de partida, y el TARIC recoge todo lo que la política comercial le suma o le resta.
\end{itemize}

\subsubsection{Modulación del arancel y regímenes aduaneros}\label{sec:modulacion}

\begin{itemize}
\item \emph{Para ajustarlo a la coyuntura}, sin cambiar el AAC:
  \begin{itemize}
  \item \concepto{contingentes arancelarios}: una cantidad de mercancía entra con un arancel reducido o nulo durante un periodo, y lo que excede paga el normal;
  \item \concepto{suspensiones arancelarias}: el arancel se suprime o se reduce, sin límite de cantidad, normalmente para insumos que no se producen en la Unión; se revisan cada 5~años;
  \item \concepto{franquicias}: exenciones permanentes por la naturaleza del envío (equipajes de viajeros, bienes personales por traslado de residencia).
  \end{itemize}
\item \emph{Regímenes aduaneros especiales}, que difieren o eximen el pago:
  \begin{itemize}
  \item \emph{depósito aduanero} y \emph{zonas francas}: la mercancía se almacena sin pagar hasta que se despacha a libre práctica;
  \item \emph{perfeccionamiento activo}: no se pagan derechos por los insumos que se transforman y se reexportan;
  \item \emph{perfeccionamiento pasivo}: al reimportar un producto que se transformó fuera, solo se paga por el valor añadido fuera;
  \item \emph{importación temporal} y \emph{tránsito}.
  \end{itemize}
\item \emph{La reforma aduanera} (2026)\nocite{CE2026ReformaAduanera}:
  \begin{itemize}
  \item un nuevo \emph{Código Aduanero de la Unión} (Reglamento 2026/2108, que se aplica de forma gradual desde septiembre de 2026) crea una \emph{Autoridad Aduanera de la UE} (en Lille) y un \emph{centro de datos aduaneros} único, obligatorio para el comercio electrónico desde 2028, y trata a las plataformas como importadores;
  \item el comercio electrónico pierde la franquicia de los envíos de menos de 150~euros: desde el 1~de julio de 2026 pagan un derecho provisional de 3~euros por artículo, y desde el 1~de noviembre de 2026, además, una tasa de gestión de 2~euros\nocite{UE2026EnviosBajoValor}.
  \end{itemize}
\end{itemize}

% --------------------------------------------------------------------------
\subsection{Las medidas de defensa comercial}\label{sec:defensa}

\begin{itemize}
\item Son los instrumentos con que la Unión protege a sus productores frente a las \emph{prácticas desleales} (el dumping y las subvenciones) y frente a las \emph{oleadas de importaciones}. Aplican los acuerdos de la OMC \vertema{3.B.33}, con algunas reglas propias más estrictas. Desde 2014, la Comisión las investiga y las impone con actos de ejecución, bajo el control de un comité de los Estados.
\end{itemize}

\subsubsection{Antidumping}\label{sec:antidumping}

\begin{itemize}
\item \emph{Definición}: hay \concepto{dumping} cuando un producto se exporta a la UE a un precio inferior a su \concepto{valor normal}: el precio en el mercado del exportador o, si no es representativo, el coste de producción más un margen razonable (Reglamento 2016/1036)\nocite{UE2016Antidumping}.
\item \emph{Requisitos} para imponer un derecho:
  \begin{itemize}
  \item un \emph{margen de dumping} de al menos el 2~\%;
  \item un \emph{perjuicio} importante (o su amenaza) para la industria de la Unión;
  \item la \emph{relación causal} entre las importaciones a bajo precio y el perjuicio;
  \item y que la medida no sea contraria al \emph{interés de la Unión} (usuarios, consumidores, importadores), una prueba que no exige la OMC.
  \end{itemize}
\item \emph{Cuánto}: el derecho no puede superar el margen de dumping y, por la \concepto{regla del derecho inferior}, se limita al margen que basta para eliminar el perjuicio. Desde la modernización de 2018 (Reglamento 2018/825) esa regla no se aplica cuando hay distorsiones en las materias primas del exportador.
\item \emph{Las economías con distorsiones significativas} (Reglamento 2017/2321): si los precios y los costes del país exportador están distorsionados por la intervención del Estado, el valor normal se construye con precios de un país representativo. Sustituyó en 2017 al trato de «economía de no mercado»; en 2025 se aplicó en los 16~casos antidumping contra China\nocite{CE2026DefensaComercial}.
\item \emph{Procedimiento}:
  \begin{itemize}
  \item la industria de la Unión denuncia (con el apoyo de al menos el 25~\% de la producción y de más de la mitad de los productores que se pronuncian), o la Comisión actúa de oficio;
  \item la Comisión abre la investigación en 45~días; puede imponer derechos provisionales a los 7~u 8~meses y definitivos en un plazo de 14~meses;
  \item las medidas duran 5~años, prorrogables tras una \emph{reconsideración por expiración}; se pueden extender a quien las elude (montaje en terceros países) y sustituir por \concepto{compromisos de precios} del exportador.
  \end{itemize}
\end{itemize}

\subsubsection{Antisubvenciones}\label{sec:antisubvenciones}

\begin{itemize}
\item \emph{Definición}: se pueden imponer \concepto{derechos compensatorios} a las importaciones que se benefician de una subvención \emph{específica} del país de origen (una ayuda pública que da una ventaja a una empresa o sector determinados), si causa perjuicio a la industria de la Unión (Reglamento 2016/1037)\nocite{UE2016Antisubvenciones}.
\item \emph{Clasificación de la OMC}: las subvenciones \emph{prohibidas} (a la exportación o al contenido local) y las \emph{recurribles} (las demás específicas, si causan efectos desfavorables). La categoría de subvenciones «no recurribles» (I+D, regiones, medio ambiente) expiró en 2000 \vertema{3.B.33}.
\item \emph{El caso de los vehículos eléctricos chinos} (Reglamento de Ejecución 2024/2754)\nocite{CE2024VehiculosElectricos}: desde el 30~de octubre de 2024 y durante 5~años, derechos compensatorios que se suman al arancel del 10~\%: el 7,8~\% para Tesla, el 17,0~\% para BYD, el 18,8~\% para Geely, el 20,7~\% para las demás empresas que cooperaron y el 35,3~\% para SAIC y las que no cooperaron. La primera investigación de oficio de este tipo; China la ha llevado a la OMC (DS630) y respondió con derechos al brandy, al porcino y a los lácteos europeos. En febrero de 2026 la Comisión aceptó el primer \emph{compromiso de precios}, de un fabricante europeo que produce en China: un precio mínimo, un límite de volumen y un compromiso de inversión en la UE.
\end{itemize}

\subsubsection{Salvaguardias}\label{sec:salvaguardias}

\begin{itemize}
\item \emph{Qué son}: protección temporal, no discriminatoria por países, frente a un aumento imprevisto de las importaciones que causa o amenaza causar un perjuicio grave (Reglamento 2015/478). No hace falta que haya una práctica desleal, pero hay que compensar a los exportadores afectados o aceptar que respondan.
\item \emph{El acero}: la UE aplicó una salvaguardia, con contingentes por países y un 25~\% fuera de ellos, desde 2018\nocite{CE2026DefensaComercial}, cuando Estados Unidos cerró su mercado y el acero se desvió hacia Europa. Expiró el 30~de junio de 2026, al llegar al máximo de 8~años de la OMC.
\item \emph{La medida que la sustituye} (Reglamento 2026/1384, desde el 1~de julio de 2026)\nocite{CE2026Salvaguardia}: ya no es una salvaguardia, sino una medida permanente basada en el art. 207~TFUE, acompañada de una renegociación de los aranceles consolidados en la OMC (art. XXVIII del GATT):
  \begin{itemize}
  \item contingentes libres de derechos de unos 18,3~millones de toneladas al año, un 47~\% menos que antes\nocite{CE2026Salvaguardia};
  \item un \emph{derecho del 50~\%} fuera de contingente;
  \item y la obligación de declarar dónde se fundió y coló el acero (\emph{melt and pour}), para evitar que el acero de un país pase por otro.
  \end{itemize}
\end{itemize}

\subsubsection{Obstáculos al comercio y cumplimiento}\label{sec:obstaculos}

\begin{itemize}
\item \emph{El Reglamento sobre obstáculos al comercio} (2015/1843): las empresas, las industrias o los Estados de la Unión pueden denunciar ante la Comisión los obstáculos de un tercer país (licencias, discriminación, falta de protección de la propiedad intelectual). Su resultado suele ser una negociación o un caso en la OMC.
\item \emph{El Reglamento de cumplimiento} (\emph{Enforcement Regulation}, 654/2014, reformado en 2021)\nocite{UE2021Ejecucion}: permite a la Unión tomar contramedidas cuando ha ganado un caso en la OMC o en un acuerdo bilateral y el otro país no cumple. Tras el bloqueo del Órgano de Apelación, también cuando el país perdedor apela «al vacío»; y cubre los servicios y la propiedad intelectual. Es la base de las contramedidas frente a Estados Unidos (\ref{sec:otras}).
\end{itemize}

\subsubsection{La defensa comercial en cifras}\label{sec:defensa-cifras}

\begin{itemize}
\item A finales de 2025 la UE tenía 232~medidas en vigor (199~un año antes): 159~antidumping, 25~antisubvenciones, 46~extensiones por elusión y 2~salvaguardias\nocite{CE2026DefensaComercial}.
\item En 2024 y 2025 abrió más investigaciones que en ningún año desde principios de siglo (gráfico~\ref{fig:defensa}): el exceso de capacidad de China, desviado de Estados Unidos, se concentra en el acero y en la química. China es, con diferencia, el país más afectado.
\end{itemize}

\begin{figure}[H]
\centering
\caption{Investigaciones de defensa comercial iniciadas y medidas impuestas por la UE, 1995-2025}\label{fig:defensa}
\grafico[0.95\textwidth]{defensa-ue}
\fuente{OMC, estadísticas de defensa comercial, fila de la Unión Europea\nocite{OMC2026DefensaUE}. La OMC cuenta por fecha de notificación: el informe anual de la Comisión da 33~investigaciones nuevas en 2024 y 32~en 2025\nocite{CE2026DefensaComercial}.}
\end{figure}

\subsubsection{Los efectos del antidumping: ¿a quién protege?}\label{sec:prusa}

\begin{itemize}
\item \emph{La justificación} es corregir la competencia desleal, pero los efectos pueden ser otros: \textcite{Prusa1992Withdrawn} observó que muchas denuncias antidumping en Estados Unidos se retiraban sin medidas y que, aun así, los precios subían.
\item \emph{La idea}: con competencia en precios, el \emph{compromiso de precios} que acepta el exportador (un precio mínimo) actúa como una coordinación de los precios. La empresa europea responde subiendo también el suyo, y las 2~ganan a costa de los consumidores y de las industrias usuarias: es una \concepto{colusión inducida} que la ley de competencia prohibiría si la pactaran las empresas.
\item \emph{La evidencia}: un derecho antidumping reduce entre un 30~y un 50~\% las importaciones del país investigado, pero las de los demás suben \parencite{Prusa2001Spread}; en la UE esa desviación es menor, porque suele investigar a varios países a la vez y extiende las medidas a quien las elude \parencite{KoningsVandenbussche2001Diversion}. Por eso la UE exige la prueba del interés de la Unión.
\item El modelo completo, con su ejemplo numérico y su gráfico, está en el anexo~\ref{anexo:prusa}.
\end{itemize}

% --------------------------------------------------------------------------
\subsection{El Sistema de Preferencias Generalizadas}\label{sec:spg}

\begin{itemize}
\item \emph{Qué es}: el \concepto{Sistema de Preferencias Generalizadas} (SPG) concede rebajas arancelarias unilaterales, sin reciprocidad, a los productos de los países en desarrollo. Nació de una propuesta de la UNCTAD (1968), la Comunidad lo aplicó desde 1971 y lo ampara la cláusula de habilitación del GATT (1979). Es a la vez política comercial y de desarrollo \vertema{3.B.42}.
\item \emph{Sus 3~regímenes}:
  \begin{itemize}
  \item el \emph{SPG general}: suprime los aranceles de los productos no sensibles y rebaja los de los sensibles, para los países de renta baja y media-baja que no tienen otro acuerdo preferencial con la UE; se retira por \emph{graduación} cuando un país deja de ser de renta media-baja o cuando un sector suyo es muy competitivo;
  \item el \emph{SPG+}: arancel cero en casi todos los productos para países vulnerables que ratifican y aplican de verdad una treintena de convenios internacionales de derechos humanos y laborales, medio ambiente y buen gobierno; en 2026 tiene 8~beneficiarios (Bolivia, Cabo Verde, Filipinas, Kirguistán, Mongolia, Pakistán, Sri Lanka y Uzbekistán)\nocite{CE2026SPGPagina};
  \item \emph{Todo menos armas} (TMA, 2001): acceso libre de aranceles y de cuotas para todo, salvo armas y municiones, de los países menos adelantados de la ONU.
  \end{itemize}
\item \emph{Su condicionalidad}: las preferencias se pueden retirar por violaciones graves de los derechos humanos. Ha ocurrido con Myanmar, Bielorrusia y Sri Lanka y, en 2020, en parte, con Camboya.
\item \emph{El nuevo reglamento} (2026/1395, de 2027 a 2036)\nocite{UE2026SPG}: sustituye al de 2012 (prorrogado hasta finales de 2026) y añade:
  \begin{itemize}
  \item la retirada de las preferencias si un país no coopera en la \emph{readmisión} de sus nacionales en situación irregular, tras un año de diálogo;
  \item el Acuerdo de París entre los convenios exigidos y una retirada más rápida en casos graves;
  \item una salvaguardia automática para el arroz;
  \item la pérdida del SPG cuando un país firma un acuerdo de libre comercio con la UE (Indonesia y Kenia salen en 2027), y una transición de 3~años para los países que dejan de ser menos adelantados (Bangladés, Laos y Nepal, en 2026).
  \end{itemize}
\item \emph{Valoración}: el SPG abre el mercado a unos 65~países, pero las reglas de origen exigentes, la erosión de las preferencias por los acuerdos bilaterales y la exclusión de los productos sensibles limitan su uso; y la condicionalidad migratoria abre el debate de si el comercio debe servir a fines ajenos al desarrollo.
\end{itemize}

% --------------------------------------------------------------------------
\subsection{Los instrumentos de seguridad\cortetitulo{} económica y sostenibilidad}\label{sec:nuevos}

\begin{itemize}
\item Desde 2019, la UE ha creado instrumentos que no protegen frente a un precio desleal, sino frente a la coerción, a la falta de reciprocidad, a los riesgos de seguridad y a la producción insostenible. Se agrupan en 2~familias (el detalle de cada uno, en la tabla~\ref{tab:nuevos})\nocite{CE2025SeguridadEconomica}:
  \begin{itemize}
  \item \emph{seguridad económica y reciprocidad}: el Instrumento Anticoerción, el de Contratación Pública Internacional, el de subvenciones extranjeras, el control de las inversiones extranjeras y el de los productos de doble uso;
  \item \emph{sostenibilidad}: el Mecanismo de Ajuste en Frontera por Carbono (MAFC), el reglamento de deforestación, la diligencia debida y la prohibición de los productos del trabajo forzoso.
  \end{itemize}
\end{itemize}

\begin{anotaciones}
\textbf{Los nuevos instrumentos, uno a uno.} Qué hace cada uno y en qué situación está en noviembre de 2026.
\begin{tablacaja}
\centering
\caption{Los nuevos instrumentos autónomos de la política comercial de la UE}\label{tab:nuevos}
\footnotesize
\begin{tabular}{>{\raggedright\arraybackslash}m{3.1cm}>{\raggedright\arraybackslash}m{5.4cm}>{\raggedright\arraybackslash}m{5.3cm}}
\toprule
\rowcolor{tceeazulpalido} \textbf{Instrumento} & \textbf{Qué hace} & \textbf{Situación en noviembre de 2026}\\
\midrule
\multicolumn{3}{l}{\textcolor{tceemorado}{\textbf{Seguridad económica y reciprocidad}}}\\
Anticoerción (Reglamento 2023/2675) & Si un país presiona a la UE o a un Estado con medidas comerciales o de inversión, permite negociar y, si no cede, responder con aranceles, restricciones de servicios, inversión o contratación pública & En vigor desde diciembre de 2023; no se ha activado nunca, aunque se planteó frente a los aranceles de Estados Unidos y frente a las restricciones chinas a las tierras raras\\
Contratación Pública Internacional (Reglamento 2022/1031) & Restringe el acceso a las licitaciones de la UE de los países que cierran las suyas & Primer uso en junio de 2025: excluye a las empresas chinas de las licitaciones de productos sanitarios de más de 5 millones de euros\\
Subvenciones extranjeras (Reglamento 2022/2560) & Investiga las subvenciones de terceros países a empresas que compran empresas europeas o licitan en la UE & Aplicable desde 2023 \vertema{3.B.39}\\
Control de inversiones (Reglamento 2026/1386) & Obliga a todos los Estados a tener un control de la inversión extranjera por seguridad, con cooperación entre ellos & Sustituye al marco de 2019 desde enero de 2028 \vertema{3.B.39}\\
Doble uso (Reglamento 2021/821) & Controla la exportación de bienes y tecnologías civiles con uso militar & Desde 2025, controles propios de la UE para las tecnologías emergentes (cuántica, semiconductores)\\
\multicolumn{3}{l}{\textcolor{tceemorado}{\textbf{Sostenibilidad}}}\\
Mecanismo de Ajuste en Frontera por Carbono (MAFC, 2023/956) & Cobra a las importaciones de cemento, acero, aluminio, fertilizantes, electricidad e hidrógeno el precio del carbono que pagan los productores europeos & Fase definitiva desde enero de 2026; certificados a la venta desde febrero de 2027 \vertema{3.B.31}\\
Deforestación (Reglamento 2023/1115) & Prohíbe vender en la UE soja, carne de vacuno, aceite de palma, cacao, café, madera y caucho de tierras deforestadas después de 2020 & Aplazado 2 veces: se aplica desde el 30 de diciembre de 2026 (las microempresas y pequeñas, desde junio de 2027)\\
Diligencia debida (Directiva 2024/1760) & Obliga a las grandes empresas a vigilar los derechos humanos y el medio ambiente en sus cadenas de suministro & Rebajada por el Ómnibus I (2026): solo empresas de más de 5.000 trabajadores y 1.500 millones de euros de facturación, desde 2029\\
Trabajo forzoso (Reglamento 2024/3015) & Prohíbe los productos hechos con trabajo forzoso & Se aplica desde diciembre de 2027\\
\bottomrule
\end{tabular}
\fuente{elaboración propia a partir de los reglamentos y de la Comisión Europea\nocite{UE2023Anticoercion,CE2025IPIProductosSanitarios,UE2026InversionesExtranjeras,UE2021DobleUso,UE2023Reglamento956,CE2025CBAMSimplificacion,CE2026Deforestacion,UE2026OmnibusI,UE2024TrabajoForzoso}.}
\end{tablacaja}
\end{anotaciones}

\begin{itemize}
\item \emph{El Instrumento Anticoerción} es el más novedoso: la Comisión examina el caso y el Consejo decide por mayoría cualificada si hay coerción; después la Comisión negocia y, si no se resuelve, adopta las contramedidas. Su fin principal es disuadir: que el agresor sepa que la respuesta llegará y será de toda la Unión.
\item \emph{El MAFC} lleva el precio europeo del carbono a la frontera para que la política climática no desplace la producción a países sin precio del carbono (\concepto{fuga de carbono}). Desde 2026, los importadores que superan 50~toneladas al año declaran las emisiones y pagan por ellas \vertema{3.B.31}.\footnote{Para respetar el trato nacional, el MAFC replica el precio del régimen europeo de comercio de derechos de emisión y descuenta el precio del carbono pagado en origen. En el tercer trimestre de 2026, el precio del certificado fue de 82,32~euros por tonelada\nocite{CE2026PrecioCBAM}. La Comisión propuso en diciembre de 2025 extenderlo a unos 180~productos transformados de acero y aluminio, que se negocia en noviembre de 2026; Rusia lo ha llevado a la OMC (DS639) \vertema{3.B.33}.}
\item \emph{Recapitulación del bloque~\ref{sec:autonoma}}: la Unión decide sola un arancel bajo pero con crestas, una defensa comercial cada vez más usada, unas preferencias cada vez más condicionadas y una caja de herramientas nueva para la seguridad y la sostenibilidad. Pero su gran palanca es la que pacta con otros.
\end{itemize}

% ==========================================================================
\section{La política comercial convencional}\label{sec:convencional}
\minutos{10}

\begin{itemize}
\item Una vez vistos los instrumentos que la Unión decide sola, pasamos a lo que pacta: en la OMC, en su red de acuerdos preferenciales y con sus 2~grandes socios sin acuerdo, Estados Unidos y China.
\item \emph{Conclusión anticipada}: ante el bloqueo de la OMC, la UE ha tejido la mayor red de acuerdos del mundo, que en 2026 cubrirá más de la mitad de su comercio; pero sus 2~primeros socios quedan fuera, y con ellos la relación se decide por el poder de negociación, no por reglas.
\end{itemize}

% --------------------------------------------------------------------------
\subsection{Los acuerdos multilaterales\cortetitulo y plurilaterales de la OMC}\label{sec:omc}

\begin{itemize}
\item \emph{Multilaterales}: la UE y sus 27~Estados están obligados por todos los acuerdos de la OMC: el GATT y los de mercancías, el Acuerdo General sobre el Comercio de Servicios y el de los Aspectos de los Derechos de Propiedad Intelectual relacionados con el Comercio (ADPIC) \vertema{3.B.33} \vertema{3.B.34}. Su política comercial se examina cada 3~años en el Mecanismo de Examen de las Políticas Comerciales, como la de las otras 3~primeras potencias comerciales.
\item \emph{Plurilaterales}: es parte del Acuerdo sobre Contratación Pública (con 22~partes, que suman 49~miembros con la UE), del de Comercio de Aeronaves Civiles y del de Tecnología de la Información (84~participantes); negoció el de Bienes Ambientales, parado desde 2016, y el de Comercio de Servicios (TiSA), suspendido desde ese mismo año; y forma parte de las iniciativas nuevas de comercio electrónico y de facilitación de las inversiones \vertema{3.B.34}.
\item \emph{Su posición}: la UE es la gran defensora del sistema multilateral. Ante el bloqueo del Órgano de Apelación impulsó en 2020 el \emph{procedimiento arbitral multipartito de apelación provisional} (MPIA), al que se han sumado 62~miembros, y propone una reforma de la OMC con plurilaterales abiertos y reglas sobre las subvenciones industriales y las empresas públicas\nocite{CE2026MPIA}.
\item \emph{Su evolución}: hasta mediados de los 2000 fue ofensiva en los temas nuevos y defensiva en la agricultura; desde el fracaso de Doha y la estrategia \emph{Europa Global} (2006), combina la defensa de la OMC con la vía bilateral.
\end{itemize}

% --------------------------------------------------------------------------
\subsection{Los acuerdos comerciales preferenciales}\label{sec:preferenciales}

\subsubsection{Tipos y generaciones}\label{sec:generaciones}

\begin{itemize}
\item \emph{La escalera de las preferencias} (figura~\ref{figu:escalera}): de menos a más integración, la UE ofrece preferencias unilaterales (SPG), acuerdos de libre comercio, acuerdos amplios y profundos con sus vecinos, uniones aduaneras y, en el Espacio Económico Europeo, el mercado interior.
\end{itemize}

\begin{infografia}[H]
\centering
\caption{La escalera de las preferencias comerciales de la UE en 2026}\label{figu:escalera}
\grafico[0.92\textwidth]{escalera-preferencias}
\fuente{elaboración propia a partir de la Comisión Europea\nocite{CE2026Negociaciones}.}
\end{infografia}

\begin{itemize}
\item \emph{4 generaciones de acuerdos}:
  \begin{itemize}
  \item \emph{vecindad y descolonización} (1960-1990): acuerdos de asociación con Grecia y Turquía (1963, con la unión aduanera de 1995), los de los países mediterráneos y las convenciones de Yaundé y de Lomé con los países de África, el Caribe y el Pacífico (ACP) \vertema{3.B.42};
  \item \emph{acuerdos de mercancías} (1990-2006): el Espacio Económico Europeo (1994), los acuerdos con Sudáfrica (2000), México (2000) y Chile (2003), centrados en los aranceles;
  \item \emph{nueva generación} (desde \emph{Europa Global}): además de los aranceles, servicios, contratación pública, propiedad intelectual e indicaciones geográficas, barreras técnicas y desarrollo sostenible. El primero fue el de Corea del Sur (2011); siguieron Colombia, Perú y Centroamérica (2013), Canadá (2017), Japón y Singapur (2019) y Vietnam (2020);
  \item \emph{acuerdos de la autonomía estratégica} (desde 2021): con capítulos de desarrollo sostenible que se pueden sancionar, comercio digital, materias primas críticas y energía. Nueva Zelanda (2024), Chile (2025), y los cerrados en 2024-2026 con el Mercosur, México, Indonesia, la India y Australia.
  \end{itemize}
\item \emph{Los capítulos de comercio y desarrollo sostenible}: tras la revisión de 2022\nocite{CE2022TSD}, los compromisos laborales (convenios fundamentales de la Organización Internacional del Trabajo) y climáticos (Acuerdo de París) se pueden hacer cumplir con sanciones comerciales como último recurso, por el procedimiento general de solución de diferencias. El primero en incluirlo fue el de Nueva Zelanda.
\end{itemize}

\subsubsection{La red de acuerdos en vigor}\label{sec:red}

\begin{itemize}
\item \emph{La mayor red del mundo}: a finales de 2025, 44~acuerdos preferenciales con 76~socios, que cubrían el 46,3~\% del comercio exterior de la UE; con el Mercosur en aplicación desde mayo de 2026, son 45~acuerdos y 80~socios\nocite{CE2026Informe}. El gráfico~\ref{fig:mapa-acuerdos} la muestra por tipos de acuerdo, y el anexo~\ref{anexo:acuerdos} la detalla por regiones.
\end{itemize}

\begin{figure}[H]
\centering
\caption{La red de acuerdos comerciales de la UE, noviembre de 2026}\label{fig:mapa-acuerdos}
\grafico[\textwidth]{mapa-acuerdos}
\fuente{elaboración propia a partir de la Comisión Europea, \emph{Negotiations and agreements}, consultada el 19/11/2026\nocite{CE2026Negociaciones}; contornos de los países de Natural Earth\nocite{NaturalEarth2026}. Cada país, con el tipo de acuerdo en aplicación (provisional o definitiva); con trama, los acuerdos cerrados o firmados y pendientes, y los que se negocian. Mónaco forma parte del territorio aduanero de la UE.}
\end{figure}

\paragraph{Europa: del mercado interior a la vecindad}

\begin{itemize}
\item \emph{El Espacio Económico Europeo} (1994) con Noruega, Islandia y Liechtenstein: el mercado interior, salvo la agricultura y la pesca, a cambio de adoptar el acervo; las diferencias las resuelve el Tribunal de la Asociación Europea de Libre Comercio (AELC).
\item \emph{Suiza}: un acuerdo de libre comercio de 1972 y unos 120~acuerdos bilaterales sectoriales. Rechazó el EEE en referéndum en 1992 (50,3~\% de noes) y abandonó en 2021 la negociación de un acuerdo marco; el paquete firmado el 2~de marzo de 2026, pendiente de ratificación, actualiza los acuerdos y les da una gobernanza común\nocite{CE2026Suiza}.
\item \emph{Las uniones aduaneras}: con Turquía (1995, productos industriales), que aplica el arancel común pero no la política comercial (los socios de la UE no están obligados a abrirle su mercado), y con Andorra y San Marino (1991), que han firmado además acuerdos de asociación.
\item \emph{La vecindad}: los \emph{acuerdos de estabilización y asociación} con los 6~países de los Balcanes Occidentales (2004-2016) y los \emph{acuerdos de libre comercio amplios y profundos} (DCFTA) con Ucrania, Moldavia y Georgia (2016), que acercan su regulación a la de la UE y preparan la adhesión\footnote{Con Ucrania, la UE suprimió en 2022 casi todos los aranceles y contingentes como apoyo frente a la invasión rusa. Al expirar esas medidas, en 2025 se revisó el art. 29~del acuerdo: contingentes más amplios, ligados a que Ucrania adopte las normas de producción de la UE (bienestar animal, plaguicidas), y una salvaguardia bilateral; se aplica desde el 29~de octubre de 2025\nocite{CE2026Ucrania}.}.
\end{itemize}

\paragraph{El Reino Unido: el Acuerdo de Comercio y Cooperación}\label{sec:reino-unido}

\begin{itemize}
\item \emph{Contexto}: tras el referéndum de 2016, el Reino Unido salió de la Unión el 31~de enero de 2020 y del mercado interior y la unión aduanera el 1~de enero de 2021 \vertema{3.B.36}. El \concepto{Acuerdo de Comercio y Cooperación} (ACC) se aplicó de forma provisional desde esa fecha y entró en vigor el 1~de mayo de 2021. Es solo UE: no lo ratificaron los Estados.
\item \emph{Contenido}:
  \begin{itemize}
  \item \emph{mercancías}: sin aranceles ni contingentes para los productos que cumplen las reglas de origen, pero con controles aduaneros y sanitarios (no hay reconocimiento mutuo);
  \item \emph{igualdad de condiciones} (\emph{level playing field}): compromisos de no rebajar las normas laborales, ambientales y de ayudas públicas, con la posibilidad de reequilibrar con aranceles;
  \item \emph{servicios}: algo más que el AGCS, pero sin pasaporte financiero ni reconocimiento automático de cualificaciones; transporte, energía, pesca (con un periodo transitorio hasta 2026), y cooperación policial y judicial.
  \end{itemize}
\item \emph{Valoración}: evitó la salida sin acuerdo, pero es un acuerdo de mínimos para quien estaba en el mercado interior. El Reino Unido sigue siendo el tercer socio comercial de la UE en bienes y el segundo en servicios. La cumbre de mayo de 2025 abrió un acercamiento: un acuerdo sanitario y fitosanitario, la conexión de los mercados de derechos de emisión y la vuelta al programa Erasmus+\nocite{CE2026ReinoUnido}. El detalle, en el anexo~\ref{anexo:reino-unido}.
\end{itemize}

\paragraph{América}

\begin{itemize}
\item \emph{Canadá, el CETA} (aplicación provisional desde 2017): el acuerdo de nueva generación más ambicioso con un país desarrollado.
  \begin{itemize}
  \item Suprime los aranceles del 99~\% de las líneas, el 98~\% desde el primer día; abre la contratación pública de las provincias y los municipios y protege 143~indicaciones geográficas.
  \item Su sistema de tribunales de inversiones solo se aplicará cuando lo ratifiquen los 27: en noviembre de 2026 faltan 9 (Chipre lo rechazó en 2020, por el queso halloumi, y el Senado francés en 2024)\nocite{CE2026CETA}.
  \item Entre 2016 y 2023 el comercio bilateral creció un 71~\%\nocite{CE2026Informe}.
  \end{itemize}
\item \emph{El Mercosur} (Argentina, Brasil, Paraguay y Uruguay), el mayor acuerdo de la UE por población y por ahorro arancelario:
  \begin{itemize}
  \item negociado desde 1999, se cerró en diciembre de 2024 y se dividió en 2: un \emph{acuerdo comercial interino}, de competencia exclusiva, y el \emph{acuerdo de asociación} completo, mixto;
  \item se firmó el 17~de enero de 2026; el acuerdo interino se aplica de forma provisional desde el 1~de mayo de 2026\nocite{CE2026MercosurProvisional}, y el de asociación espera al dictamen 1/26 que el Parlamento Europeo pidió al Tribunal de Justicia en enero\nocite{PE2026Mercosur} \vertema{3.B.35};
  \item el Mercosur liberaliza el 91~\% de las exportaciones de la UE (con aranceles de partida del 35~\% en los automóviles) y la UE el 92~\% de las suyas, con contingentes para los productos agrícolas sensibles (99.000~toneladas de carne de vacuno al 7,5~\%), una salvaguardia bilateral reforzada y la protección de 344~indicaciones geográficas europeas\nocite{CE2026MercosurAcuerdo}. La UE ya le vendía 57.000~millones de euros en bienes en 2024.
  \end{itemize}
\item \emph{México}: el Acuerdo Global de 2000, de primera generación, centrado en las mercancías. Su modernización, cerrada en enero de 2025 y firmada el 22~de mayo de 2026, lo convierte en uno de nueva generación (\ref{sec:negociacion})\nocite{CE2026Mexico}.
\item \emph{Chile}: el acuerdo de asociación de 2003 se ha modernizado con un \emph{Acuerdo Marco Avanzado} (mixto) y un \emph{Acuerdo Comercial Interino} (solo UE), en vigor desde el 1~de febrero de 2025, que suprime casi todos los aranceles restantes e incluye un capítulo de materias primas (litio, cobre) e hidrógeno\nocite{CE2026Chile}.
\item \emph{Colombia, Perú y Ecuador} (acuerdo comercial multipartes, 2013; Ecuador desde 2017) y \emph{Centroamérica} (el pilar comercial del acuerdo de asociación de 2013 con Costa Rica, El Salvador, Guatemala, Honduras, Nicaragua y Panamá): de nueva generación, con servicios, contratación pública y desarrollo sostenible.
\end{itemize}

\paragraph{Asia y Oceanía}

\begin{itemize}
\item \emph{Corea del Sur} (aplicación provisional en 2011, en vigor en 2015): el primero de nueva generación; suprimió casi todos los aranceles de forma escalonada y abordó las barreras técnicas del automóvil y la electrónica. En junio de 2026 se le sumó un acuerdo de comercio digital\nocite{CE2026Corea}.
\item \emph{Japón} (acuerdo de asociación económica, 2019): el mayor por volumen hasta el Mercosur (la UE le vende cerca de 70.000~millones de euros en bienes y 28.000~en servicios al año). Elimina los aranceles japoneses a muchos quesos, vinos y carnes, protege más de 200~indicaciones geográficas y abre la contratación pública de 48~grandes ciudades; la UE abre su mercado del automóvil en 7~años. En 2023 se le añadió un acuerdo sobre flujos de datos\nocite{CE2026Japon}.
\item \emph{Singapur} (2019) y \emph{Vietnam} (2020), los primeros con la ASEAN, a la espera de un acuerdo regional. El de Vietnam, el más ambicioso con un país en desarrollo, elimina casi el 99~\% de los aranceles, los vietnamitas en 10~años. Sus acuerdos de protección de inversiones, mixtos, siguen pendientes de ratificación; con Singapur está en vigor desde febrero de 2026 el primer acuerdo de comercio digital de la UE\nocite{CE2026Singapur,CE2026Vietnam}.
\item \emph{Nueva Zelanda} (2024): el primero con un capítulo de desarrollo sostenible que se puede hacer cumplir con sanciones (Acuerdo de París y convenios fundamentales de la Organización Internacional del Trabajo); ahorra a las empresas europeas unos 140~millones de euros de derechos al año\nocite{CE2026NuevaZelanda}.
\end{itemize}

\paragraph{África, el Caribe, el Pacífico y el Mediterráneo}

\begin{itemize}
\item \emph{Los acuerdos de asociación económica} (AAE) con los países de África, el Caribe y el Pacífico (ACP): el Cariforum (2008), la SADC (2016), África Oriental y Austral (2012), Camerún (2014), Costa de Marfil y Ghana (2016), Kenia (2024) y 4~países del Pacífico. Son asimétricos: la UE abre todo su mercado desde el primer día y sus socios, solo una parte y de forma gradual. Los regionales de África Occidental y de la Comunidad de África Oriental siguen pendientes \vertema{3.B.42}.
\item \emph{Los acuerdos euromediterráneos de asociación} (1997-2006) con Marruecos, Argelia, Túnez, Egipto, Jordania, Líbano, Israel y la Autoridad Palestina: libre comercio industrial y concesiones agrícolas recíprocas; su modernización hacia acuerdos amplios y profundos está parada.
\end{itemize}
\subsubsection{¿Qué efecto tienen los acuerdos?}\label{sec:efectos}

\begin{itemize}
\item \emph{La teoría}: un acuerdo preferencial crea comercio, pero también puede desviarlo de terceros más eficientes \vertema{3.B.10}. Con aranceles exteriores bajos, como los de la UE, el riesgo de desviación es pequeño, y los acuerdos de nueva generación añaden ganancias que el arancel no recoge: menos barreras técnicas, acceso a la contratación pública y seguridad jurídica.
\item \emph{Lo que muestran los datos}:
  \begin{itemize}
  \item a finales de 2025, los 44~acuerdos preferenciales con 76~socios cubrían el 46,3~\% del comercio exterior de la UE, 2,39~billones de euros; con los acuerdos cerrados en 2025 y 2026, la cobertura subirá al 53,3~\%\nocite{CE2026Informe};
  \item en 2025, el comercio con los socios preferenciales creció un 3,1~\%, frente al 1,4~\% con los demás;
  \item con modelos de gravedad, \textcite{BaierBergstrand2007} estiman que un acuerdo de libre comercio duplica, de media, el comercio entre sus miembros en unos 10~años.
  \end{itemize}
\end{itemize}

% --------------------------------------------------------------------------
\subsection{Acuerdos en negociación\cortetitulo y pendientes de ratificación}\label{sec:negociacion}

\begin{itemize}
\item 2025 y 2026 han sido los años de más acuerdos cerrados de la historia de la PCC: ante los aranceles de Estados Unidos, la UE ha acelerado la diversificación (tabla~\ref{tab:negociacion}).
\end{itemize}

\begin{table}[H]
\centering
\caption{Los acuerdos de la UE cerrados o en negociación, noviembre de 2026}\label{tab:negociacion}
\footnotesize
\begin{tabular}{>{\raggedright\arraybackslash}m{2.6cm}>{\raggedright\arraybackslash}m{6.2cm}>{\raggedright\arraybackslash}m{5.8cm}}
\toprule
\rowcolor{tceeazulpalido} \textbf{Socio} & \textbf{Situación} & \textbf{Contenido destacado}\\
\midrule
Mercosur & Firmado en enero de 2026; acuerdo interino en aplicación provisional desde mayo; el de asociación, pendiente del Tribunal de Justicia & Libera el 91 y el 92~\% del comercio; contingentes agrícolas\\
México & Acuerdo global modernizado cerrado en enero de 2025 y firmado el 22 de mayo de 2026; acuerdo interino aprobado por la UE en julio, pendiente de México & Casi todos los aranceles a cero, protección de inversiones, contratación pública\\
India & Negociación relanzada en 2022 y cerrada el 27 de enero de 2026; en septiembre la Comisión propuso su firma & Rebajas en más del 96~\% de las exportaciones de la UE (vino, aceite, automóviles); inversiones e indicaciones geográficas, aparte\\
Indonesia & Acuerdo de asociación económica integral cerrado en septiembre de 2025; propuestas de firma al Consejo en junio de 2026 & Indonesia sale del SPG en 2027; aceite de palma y níquel\\
Australia & Negociación relanzada en 2025 y cerrada el 24 de marzo de 2026 & Productos agrícolas y materias primas críticas\\
Suiza & Paquete firmado el 2 de marzo de 2026 & Actualiza los acuerdos bilaterales y su gobernanza\\
Tailandia, Filipinas, Malasia & En negociación & Hacia una red con la ASEAN\\
Emiratos Árabes Unidos & En negociación desde 2025; con el resto del Consejo de Cooperación del Golfo, suspendida & Energía e inversiones\\
China (inversiones) & Acuerdo Global de Inversiones cerrado en 2020 y congelado por el Parlamento en 2021 & Acceso al mercado y empresas públicas\\
\bottomrule
\end{tabular}
\fuente{elaboración propia a partir de la Comisión Europea\nocite{CE2026Negociaciones,CE2026India,CE2026Mexico,CE2026Indonesia,CE2026Australia,CE2026Informe}.}
\end{table}

% --------------------------------------------------------------------------
\subsection{Otras relaciones comerciales:\cortetitulo Estados Unidos, China y Rusia}\label{sec:otras}

\subsubsection{Estados Unidos}\label{sec:eeuu}

\begin{itemize}
\item \emph{El primer socio, sin acuerdo}: el Acuerdo Transatlántico de Comercio e Inversiones (TTIP), negociado entre 2013 y 2016, se paró con la llegada de \autor{Trump} y el Consejo dio por obsoletas sus directrices en 2019.
\item \emph{La primera administración \autor{Trump}} (2018-2020): aranceles al acero (25~\%) y al aluminio (10~\%), a los que la UE respondió con contramedidas, y la disputa de \emph{Airbus}-\emph{Boeing}, con aranceles de los 2~lados autorizados por la OMC y una tregua en 2021 \vertema{3.B.8}.
\item \emph{2025}\nocite{UE2025DeclaracionEEUU,JonesDay2026IEEPA}:
  \begin{itemize}
  \item Estados Unidos subió los aranceles al acero y al aluminio al 50~\%, puso un 25~\% a los automóviles y anunció un arancel «recíproco» del 20~\% para la UE, rebajado enseguida a un 10~\% general;
  \item la UE preparó contramedidas por 93.000~millones de euros, pero eligió negociar: en el pacto de Turnberry (27~de julio) y la declaración conjunta del 21~de agosto, Estados Unidos aplica a la UE un \emph{techo del 15~\%} que incluye el arancel NMF y cubre también los automóviles, los productos farmacéuticos y los semiconductores, y la UE suprime sus aranceles a los productos industriales estadounidenses y abre contingentes para algunos agroalimentarios;
  \item además, la UE manifestó la intención de comprar energía por 750.000~millones de dólares y de que sus empresas inviertan 600.000~millones en Estados Unidos.
  \end{itemize}
\item \emph{2026}:
  \begin{itemize}
  \item las rebajas de la UE se aplican, con 2~reglamentos, desde el 1~de julio de 2026 hasta finales de 2029, con una cláusula para suspenderlas si los aranceles de Estados Unidos al acero y al aluminio siguen por encima del 15~\% después de 2026\nocite{CE2026EEUU};
  \item el 20~de febrero, el Tribunal Supremo de Estados Unidos anuló los aranceles basados en la ley de emergencia (IEEPA); el gobierno los sustituyó por un recargo del 10~\% (sección~122) y, desde julio, por aranceles de la sección~301, que a la UE se le aplican descontando el arancel NMF para respetar el techo del 15~\% \vertema{3.B.33};
  \item en julio, la UE suspendió sin fecha sus contramedidas, que puede reactivar en cualquier momento\nocite{CE2026Contramedidas}. El acero, el aluminio y el cobre siguen pagando el 50~\%.
  \end{itemize}
\end{itemize}

\subsubsection{Aranceles óptimos, represalias y negociación:\cortetitulo{} el acuerdo con Estados Unidos}\label{sec:negociacion-modelo}

\begin{itemize}
\item \emph{El punto de partida}: un país grande puede ganar con su arancel óptimo, pero si el otro responde, la guerra arancelaria deja a los 2~peor que el libre comercio: es el equilibrio de \autor{Nash}\footnote{\autor{John Nash} recibió el Premio Nobel de Economía en 1994, con \autor{John Harsanyi} y \autor{Reinhard Selten}, «por su análisis pionero de los equilibrios en la teoría de los juegos no cooperativos»\nocite{Nobel1994Nash}.} en aranceles de \textcite{Johnson1953Optimum} \vertema{3.B.7} \vertema{3.B.8}. Las reglas de la OMC están para evitarla \vertema{3.B.33}.
\item \emph{La idea}: cuando las reglas no se respetan, los países negocian, y el reparto lo decide lo que cada uno pierde si no hay acuerdo, su \concepto{punto de desacuerdo}. Es la solución de negociación de \textcite{Nash1950Bargaining}, que \textcite{GrossmanHelpman1995TradeWars} aplican a las guerras comerciales: gana más quien menos pierde con la ruptura.
\item \emph{Lo que explica}:
  \begin{itemize}
  \item \emph{el acuerdo asimétrico de 2025}: la UE aceptó un 15~\% frente a un arancel cero propio porque, en una guerra comercial, perdía más que Estados Unidos (exporta mucho más de lo que le compra en bienes) y porque la negociación estaba unida a la seguridad de Europa;
  \item \emph{el valor de la amenaza}: las contramedidas preparadas y el Instrumento Anticoerción no sirven tanto para usarse como para empeorar el punto de desacuerdo del otro y mejorar el reparto. Es «donde las dan, las toman» como disuasión: la sola voz de \autor{Meunier} (\ref{sec:meunier}) solo pesa si detrás hay capacidad de responder;
  \item \emph{los datos}: el Reino Unido, con un comercio de bienes más equilibrado con Estados Unidos, obtuvo un 10~\%; la UE y Japón, con grandes superávits, un 15~\%.
  \end{itemize}
\item \emph{Su límite}: explica el reparto, pero no que el acuerdo sea estable (es una declaración política, no un tratado), y a largo plazo el mejor punto de desacuerdo de la UE es tener alternativas, que es lo que construye su red de acuerdos. El modelo completo, con su gráfico, está en el anexo~\ref{anexo:nash}.
\end{itemize}

\subsubsection{China}\label{sec:china}

\begin{itemize}
\item \emph{Socio, competidor y rival sistémico} (así la define la UE desde 2019): es el primer proveedor de la UE y el segundo socio comercial, con el mayor déficit bilateral (375.000~millones de euros en 2025).
\item \emph{De la apertura a la reducción de riesgos} (\emph{de-risking}): el Acuerdo Global de Inversiones, cerrado en diciembre de 2020, quedó congelado por el Parlamento Europeo en mayo de 2021 tras las sanciones cruzadas. Desde entonces, la UE usa la defensa comercial (los vehículos eléctricos), los instrumentos nuevos (la contratación pública de productos sanitarios, las subvenciones extranjeras) y la diversificación de los insumos críticos.
\item \emph{La respuesta china}: derechos antidumping al brandy (2025), al porcino (2025) y antisubvenciones a los lácteos (2026), y controles a la exportación de tierras raras e imanes desde abril de 2025, ampliados en octubre y suspendidos en parte después\nocite{EPRS2025TierrasRaras}. El exceso de capacidad chino, desviado del mercado estadounidense, es hoy el principal motivo de la defensa comercial europea.
\end{itemize}

\subsubsection{Rusia y los países ACP}\label{sec:rusia}

\begin{itemize}
\item \emph{Rusia}: el Acuerdo de Colaboración y Cooperación de 1997 quedó vacío con las sanciones. Desde la invasión de Ucrania en 2022, la UE ha aprobado más de 20~paquetes de sanciones (el 21.º, en julio de 2026)\nocite{CONS2026Sanciones21}, ha puesto aranceles prohibitivos a sus fertilizantes y productos agrícolas (2025) y dejará de importar su gas: el gas natural licuado desde 2027 y el de gasoducto a finales de ese año (Reglamento 2026/261)\nocite{UE2026Gas}. Rusia, que era el quinto socio comercial de la UE en 2021, cayó en 2025 al puesto 19\nocite{Eurostat2026Socios}.
\item \emph{Los países ACP}: sus relaciones con la UE se rigen por el Acuerdo de Samoa (2023), sucesor del de Cotonú, y su vertiente comercial son el SPG y los acuerdos de asociación económica; se estudian con las relaciones exteriores y la cooperación al desarrollo de la Unión \vertema{3.B.42}.
\end{itemize}

% --------------------------------------------------------------------------
\subsection{Retos de futuro}\label{sec:retos}

\begin{itemize}
\item Una vez vista la articulación de la política, cerramos con los retos que tiene por delante. Son 5, y todos giran en torno a cómo seguir abierta en un mundo menos cooperativo:
  \begin{itemize}
  \item \emph{Sostener el multilateralismo sin la OMC de antes}: con el Órgano de Apelación bloqueado y las grandes potencias fuera de las reglas, la UE tiene que decidir hasta dónde apuesta por los acuerdos plurilaterales y por el MPIA, y cómo reforma la OMC para que trate las subvenciones industriales y las empresas públicas \vertema{3.B.33}.
  \item \emph{Dar estabilidad a la relación con Estados Unidos}: el pacto de 2025 es una declaración política, sin solución de diferencias, que depende de las bases legales estadounidenses y de que el acero y el aluminio vuelvan al techo del 15~\% \vertema{3.B.33}.
  \item \emph{Responder al exceso de capacidad de China sin cerrarse}: la defensa comercial y los nuevos instrumentos se usan cada vez más, y el riesgo es que la seguridad económica se convierta en proteccionismo; el Consejo Europeo pidió en octubre de 2026 reforzar la caja de herramientas\nocite{CEU2026ConclusionesOctubre}.
  \item \emph{Ratificar a tiempo}: la politización de los acuerdos (el CETA, el Mercosur ante el Tribunal de Justicia) alarga los plazos; la arquitectura de 2018 y la mayoría cualificada son la respuesta, y la India, Indonesia y Australia la pondrán a prueba.
  \item \emph{Conciliar sostenibilidad y competitividad}: el MAFC, la deforestación y la diligencia debida exportan las normas europeas, pero sus socios los ven como barreras y las empresas europeas, como cargas; los aplazamientos y el Ómnibus de 2026 muestran la tensión \vertema{3.B.31}.
  \end{itemize}
\item \emph{Recapitulación del bloque~\ref{sec:convencional}}: la UE defiende la OMC, pero su gran instrumento es la red de acuerdos, que cubrirá más de la mitad de su comercio; con Estados Unidos negocia bajo amenaza, con China se protege, y su reto es seguir abierta sin quedar indefensa.
\end{itemize}

% ==========================================================================
\section*{Conclusión}
\minutos{1}

\begin{itemize}
\item \textbf{\emph{Recapitulación}}: En conclusión,
  \begin{itemize}
  \item La política comercial común es la consecuencia de la unión aduanera y la fuente del poder de negociación de la Unión; su alcance lo han ido fijando los Tratados y el Tribunal de Justicia, y su eficacia depende de cómo se repartan el voto y la delegación.
  \item La Unión decide sola un arancel bajo, una defensa comercial que crece con el exceso de capacidad mundial, unas preferencias para los países en desarrollo y una caja de herramientas nueva para la seguridad y la sostenibilidad.
  \item Y pacta la mayor red de acuerdos del mundo, mientras negocia con Estados Unidos bajo amenaza y se protege de China.
  \end{itemize}
\item \textbf{\emph{Relevancia}}:
  \begin{itemize}
  \item Para España, que exporta fuera de la Unión casi 4~de cada 10~euros de sus ventas de bienes, la PCC decide su acceso a los mercados (el aceite, el vino, el automóvil) y su defensa frente a la competencia desleal \vertema{4.A.22} \vertema{4.A.24}.
  \end{itemize}
\item \textbf{\emph{Extensiones y relación con otras partes del temario}}:
  \begin{itemize}
  \item Las reglas multilaterales \vertema{3.B.33} \vertema{3.B.34}; la teoría de la integración \vertema{3.B.10} y los bloques fuera de la UE \vertema{3.B.35}; el mercado interior y el control de inversiones \vertema{3.B.39}; el cambio climático y el MAFC \vertema{3.B.31}; las relaciones exteriores y la cooperación al desarrollo \vertema{3.B.42}; y la política española de internacionalización \vertema{4.A.24} \vertema{4.A.25}.
  \end{itemize}
\item \textbf{\emph{Opinión}}:
  \begin{itemize}
  \item La seguridad económica es necesaria, pero corre un riesgo: que cada instrumento nuevo acabe protegiendo a quien más presiona en lugar de responder a la coerción.

  \item Y la red de acuerdos, la mejor defensa de la Unión, necesita que se cierren por mayoría cualificada.
  \end{itemize}
\item \textbf{\emph{Idea final}}:
  \begin{itemize}
  \item Volvamos al refrán: donde las dan, las toman. En 2025 la Unión no devolvió el golpe, pero el modelo de la negociación enseña que la amenaza solo pesa si es creíble y si se tienen alternativas. Una potencia comercial que vive de la apertura no se defiende cerrándose, sino teniendo instrumentos para responder y muchos socios con los que comerciar.
  \end{itemize}
\end{itemize}

\clearpage
\appendix

% ==========================================================================
\section{Anexo: la economía política\cortetitulo{} de la sola voz (\autor{Meunier}, 2005)}\label{anexo:meunier}

\begin{anotaciones}
\textbf{Notación.} $x$ es el grado de apertura del mercado de la UE, $\mathit{SQ}$ el statu quo, $x_i$ el punto ideal del Estado $i$, $x^{O}$ el del socio negociador, $x_p$ el del Estado pivote y $x^*$ el resultado.
\end{anotaciones}

\paragraph{a) Idea}

\begin{itemize}
\item Que la UE hable con una sola voz no garantiza, por sí solo, que negocie mejor. \textcite{Meunier2005Trading} muestra que el resultado depende de 2~variables institucionales y de una de contexto:
  \begin{itemize}
  \item la \emph{regla de voto} en el Consejo: unanimidad o mayoría cualificada;
  \item la \emph{delegación} en la Comisión: un mandato restrictivo y vigilado o uno amplio;
  \item y el \emph{contexto}: si el socio quiere cambiar la política de la UE (caso \emph{conservador} o defensivo, como la agricultura) o es la UE la que quiere cambiar la del socio (caso \emph{reformista} u ofensivo, como la contratación pública).
  \end{itemize}
\item Es la aplicación a la Unión de la idea de \autor{Thomas Schelling} (1960): en una negociación, tener las manos atadas puede ser una ventaja\nocite{Schelling1960Strategy}, y de los juegos a 2~niveles de \textcite{Putnam1988Diplomacy}: lo que se pacta fuera tiene que poder ratificarse dentro.
\end{itemize}

\paragraph{b) Modelo}

\subparagraph{i) Supuestos}

\begin{itemize}
\item Un solo eje de política: el grado de apertura $x$ del mercado de la UE (o de su socio, en el caso reformista), con el \emph{statu quo} en $\mathit{SQ}=0$.
\item $n$ Estados miembros con puntos ideales $x_1<x_2<\dots<x_n$ y preferencias simétricas (por ejemplo, cuadráticas) en torno a ellos: el Estado $i$ acepta cualquier acuerdo que le deje al menos tan bien como el \emph{statu quo}, es decir, cualquier $x\in[0,\,2x_i]$.
\item Un socio negociador $O$ que quiere la mayor apertura posible ($x^{O}$ muy por encima de los $x_i$) y que puede hacer una oferta del tipo «lo tomas o lo dejas».
\item El acuerdo necesita la aprobación del Consejo: por unanimidad o por mayoría cualificada. El \concepto{Estado pivote} es el último cuyo voto hace falta para aprobarlo.
\end{itemize}

\subparagraph{ii) Desarrollo}

\begin{itemize}
\item \emph{Paso 1: el conjunto ganador}.
  \begin{itemize}
  \item El \concepto{conjunto ganador} (\emph{win-set}) de la UE son los acuerdos que el Consejo aprobaría frente al \emph{statu quo}. Lo fija el pivote: el Estado $p$ acepta hasta $2x_p$, y todos los Estados más abiertos que él, también.
  \end{itemize}
\item \emph{Paso 2: el resultado}.
  \begin{itemize}
  \item El socio empuja hasta el borde del conjunto ganador, o hasta su punto ideal si está dentro:
  \funcionclave[eq:meunier]{x^*=\min\big\{x^{O},\;2x_p\big\},\qquad p=\begin{cases}1 & \text{unanimidad}\\ q & \text{mayoría cualificada}\end{cases}}
  \end{itemize}
\item \emph{Paso 3: comparar las reglas} (gráfico~\ref{fig:meunier}).
  \begin{itemize}
  \item Con unanimidad decide el Estado más conservador, $x_1$, y el resultado es $U=2x_1$.
  \item Con mayoría cualificada decide un Estado intermedio $x_q>x_1$, y el resultado es $M=2x_q>U$: más apertura y un acuerdo mejor para el socio.
  \end{itemize}
\item \emph{Paso 4: la delegación}.
  \begin{itemize}
  \item Un mandato restrictivo y vigilado (el Comité de Política Comercial, que revisa cada paso) hace creíble que la Comisión no puede ir más allá de $2x_p$: son las \emph{manos atadas}.
  \item Un mandato amplio da a la Comisión margen para cerrar acuerdos, pero el socio sabe que puede sacarle más.
  \end{itemize}
\end{itemize}

\begin{figure}[H]
\centering
\caption{La economía política de la sola voz en el caso conservador}\label{fig:meunier}
\grafico[0.92\textwidth]{meunier}
\fuente{elaboración propia a partir de \textcite[cap.~2]{Meunier2005Trading}. Puntos ideales ilustrativos de 5~Estados; con unanimidad decide $x_1$ y con mayoría cualificada (3~de 5), $x_3$.}
\end{figure}

\subparagraph{iii) Implicaciones}

\begin{itemize}
\item \emph{En el caso conservador, la unanimidad y la delegación restrictiva hacen de la UE un negociador duro}: el socio no obtiene más de lo que concede el Estado más reacio, y la Unión lo hace valer con el peso de todo su mercado. Los Estados más proteccionistas son los que más ganan con la sola voz.
\item \emph{En el caso reformista, las reglas internas importan menos}: el resultado depende de lo que acepte el socio; pero con mayoría cualificada la UE consigue lanzar la ofensiva y cerrar el acuerdo, y ningún Estado aislado tendría peso para hacerlo.
\item \emph{Sin competencia exclusiva, el socio divide}: puede pactar con los Estados más abiertos uno a uno, como hizo Estados Unidos con los acuerdos bilaterales de «cielos abiertos» en la aviación en los años 90.
\item \emph{Es un dilema de diseño}: más mayoría y más delegación dan eficacia y acuerdos; más unanimidad y más control dan dureza y legitimidad.
\end{itemize}

\paragraph{c) Extensiones}

\begin{itemize}
\item \emph{La politización}: desde el TTIP y el CETA, los parlamentos nacionales y regionales y la opinión pública son un nivel más de ratificación; con \autor{Putnam}, un conjunto ganador interno más pequeño da más fuerza ante el socio, pero también más riesgo de que no haya acuerdo.
\item \emph{El Parlamento Europeo como veto}: desde Lisboa tiene que aprobar todos los acuerdos, y lo ha usado (rechazó el acuerdo antifalsificación ACTA en 2012 y en 2026 retrasó el Mercosur pidiendo un dictamen al Tribunal de Justicia).
\end{itemize}

\paragraph{d) Evidencia empírica}

\begin{itemize}
\item \textcite{Meunier2005Trading} contrasta el modelo con 4~negociaciones con Estados Unidos: en la agricultura de la Ronda Kennedy (unanimidad, caso conservador) Estados Unidos acabó aceptando la posición de mínimos de la Comunidad; en la Ronda Uruguay, el preacuerdo de Blair House (1992) se renegoció cuando Francia reafirmó su veto; en la contratación pública (1990-1994, caso reformista con mayoría) la Comunidad abrió en parte el mercado estadounidense, y en los cielos abiertos, sin competencia común, Estados Unidos dividió a los europeos.
\item Lo mismo se ve hoy: el acuerdo con el Mercosur salió adelante en el Consejo por mayoría cualificada, con Francia en contra, porque se dividió en un acuerdo comercial de competencia exclusiva y un acuerdo de asociación mixto\nocite{Consejo2026Mercosur}.
\end{itemize}

\paragraph{e) Valoración}

\begin{itemize}
\item Explica bien por qué la UE es dura en lo defensivo y lenta en lo ofensivo, y por qué la regla de voto y la división de los acuerdos importan tanto.
\item Pero reduce la negociación a un eje y supone que el socio no tiene sus propias restricciones internas; deja fuera el poder de cada parte para \emph{amenazar} si no hay acuerdo, que es lo que decidió la negociación con Estados Unidos en 2025 (anexo~\ref{anexo:nash}).
\end{itemize}

% ==========================================================================
\section{Anexo: el reparto de competencias\cortetitulo{} de la Unión}\label{anexo:competencias}

\begin{itemize}
\item \emph{La regla de partida es la atribución}: la Unión solo tiene las competencias que le dan los Tratados, y las demás siguen siendo de los Estados (arts.~4.1~y 5~TUE). El art.~2~TFUE las ordena en categorías según quién legisla, y los arts.~3~a 6~enumeran los ámbitos de cada una (figura~\ref{figu:competencias-ue}). Las ejerce con \emph{proporcionalidad} (no más de lo necesario) y, fuera de las exclusivas, con \emph{subsidiariedad} (solo si el objetivo se alcanza mejor en común).
\item \emph{Exclusivas y compartidas}: en las \concepto{exclusivas} (art.~3) solo legisla la Unión; los Estados, solo si ella los faculta o para aplicar sus actos. En las \concepto{compartidas} (art.~4), la categoría residual, legislan las 2, pero lo que regula un acto de la Unión deja de estar disponible para los Estados: es la \emph{preferencia} u ocupación del campo (art.~2.2~y Protocolo n.º~25), que solo alcanza a lo que el acto regula. En investigación, espacio, cooperación para el desarrollo y ayuda humanitaria son \emph{paralelas}: la acción de la Unión no impide la de los Estados (art.~4.3~y 4.4).
\item \emph{Coordinación y apoyo}: en las políticas económicas y de empleo (art.~5) la política sigue siendo de los Estados y la Unión la coordina con orientaciones del Consejo; en los ámbitos del art.~6 (salud, industria, cultura, turismo, educación, protección civil, cooperación administrativa) la Unión solo apoya, coordina o complementa, sin armonizar las leyes nacionales. La PESC tiene un régimen propio, intergubernamental, con unanimidad y sin actos legislativos (art.~24~TUE).
\item \emph{Dónde está este tema}: la política comercial común (art.~3.1.e; arts.~206~y 207~TFUE) y la unión aduanera (art.~3.1.a) son exclusivas, y también lo es celebrar acuerdos internacionales cuando lo prevé un acto de la Unión, es necesario para ejercer una competencia interna o puede afectar a normas comunes (art.~3.2).
\item \emph{Qué se deduce}: la Unión habla con una sola voz y los Estados no pueden firmar acuerdos comerciales propios. Pero cuando un acuerdo entra en terrenos compartidos (las inversiones de cartera o el arbitraje entre inversores y Estados, \emph{Dictamen 2/15}), es mixto y lo ratifican también los Estados.
\end{itemize}

\begin{infografia}[H]
\centering
\caption{El reparto de competencias de la Unión (arts.~2~a 6~TFUE): la política comercial común}\label{figu:competencias-ue}
\grafico[\textwidth]{competencias-ue}
\fuente{elaboración propia a partir de los arts.~2~a 6~del TFUE y de los arts.~4, 5~y 24~del TUE\nocite{UE2016TFUE}.}
\end{infografia}

\clearpage
% ==========================================================================
\section{Anexo: el antidumping\cortetitulo{} como colusión inducida (\autor{Prusa}, 1992)}\label{anexo:prusa}

\begin{anotaciones}
\textbf{Notación.} $p_D$ y $p_F$ son los precios de la empresa de la UE (doméstica) y de la extranjera, $c$ el coste marginal, $b$ el grado de sustituibilidad, $\bar p$ el precio mínimo y $\pi_i$ los beneficios.
\end{anotaciones}

\paragraph{a) Idea}

\begin{itemize}
\item La defensa comercial se justifica como una corrección de la competencia desleal, pero sus efectos pueden ser otros: \textcite{Prusa1992Withdrawn} mostró que muchas denuncias antidumping en Estados Unidos se retiraban sin medidas y que, aun así, los precios subían: la mera amenaza de un derecho servía para que las empresas coordinaran sus precios.
\item En la UE, el instrumento equivalente es el \emph{compromiso de precios}, que fija un precio mínimo al exportador. ¿A quién beneficia?
\end{itemize}

\paragraph{b) Modelo}

\subparagraph{i) Supuestos}

\begin{itemize}
\item 2~empresas, una de la UE ($D$) y una extranjera ($F$), compiten en precios (\autor{Bertrand}) en el mercado de la UE con productos diferenciados.
\item Demandas lineales: $q_i=a-p_i+b\,p_j$, con $0<b<1$ (sustitutivos imperfectos), y el mismo coste marginal constante $c$.
\item La investigación antidumping termina con un compromiso: la empresa extranjera no venderá por debajo de un precio mínimo $\bar p$; la de la UE elige su precio después.
\end{itemize}

\subparagraph{ii) Desarrollo}

\begin{itemize}
\item \emph{Paso 1: las funciones de reacción}.
  \begin{itemize}
  \item Cada empresa maximiza $\pi_i=(p_i-c)(a-p_i+b\,p_j)$, y su mejor respuesta sube con el precio de la rival (los precios son \emph{complementos estratégicos}):
  \funcionclave[eq:reaccion]{p_i=R_i(p_j)=\frac{a+c+b\,p_j}{2}}
  \end{itemize}
\item \emph{Paso 2: el equilibrio sin antidumping}.
  \begin{itemize}
  \item En el equilibrio de \autor{Nash} los 2~precios son iguales: $p^N=\dfrac{a+c}{2-b}$.
  \item Con $a=10$, $c=2$ y $b=0{,}8$: $p^N=10$ y cada empresa gana $\pi^N=64$ (punto $N$ del gráfico~\ref{fig:prusa}).
  \end{itemize}
\item \emph{Paso 3: el compromiso de precios}.
  \begin{itemize}
  \item La extranjera vende a $\bar p>p^N$ y la de la UE responde subiendo también su precio, a $p_D=R_D(\bar p)$.
  \item Con $\bar p=12$: $p_D=10{,}8$; la empresa de la UE gana $77{,}4$ y la extranjera $66{,}4$, las 2~más que en $N$ (punto $C$).
  \end{itemize}
\item \emph{Paso 4: cuándo ganan las 2}.
  \begin{itemize}
  \item La de la UE gana siempre, porque su rival se ha atado a un precio más alto.
  \item La extranjera gana mientras $\bar p$ no sea demasiado alto (aquí, entre 10~y 13,8): pierde cuota, pero vende más caro, como un líder de \autor{Stackelberg} en precios. El punto $C$ queda dentro de la lente entre las curvas de isobeneficio que pasan por $N$.
  \end{itemize}
\end{itemize}

\begin{figure}[H]
\centering
\caption{El compromiso de precios como colusión inducida}\label{fig:prusa}
\grafico[0.7\textwidth]{prusa}
\fuente{elaboración propia a partir de \textcite{Prusa1992Withdrawn}, con el ejemplo del texto ($a=10$, $c=2$, $b=0{,}8$, $\bar p=12$). La zona sombreada son los precios con los que las 2~empresas ganan más que en $N$.}
\end{figure}

\subparagraph{iii) Implicaciones}

\begin{itemize}
\item \emph{El antidumping puede facilitar la colusión}: el precio mínimo hace de coordinación que la ley de competencia prohibiría si la acordaran las empresas entre sí.
\item \emph{Pierden los consumidores y las industrias usuarias}: los 2~precios suben, y por eso la UE exige la prueba del interés de la Unión.
\item \emph{Una medida contra un país desvía el comercio}: si solo se grava a un exportador, los demás ocupan su lugar.
\end{itemize}

\paragraph{c) Extensiones}

\begin{itemize}
\item \emph{Con cantidades} (\autor{Cournot}) el derecho no tiene este efecto colusivo, porque las cantidades son sustitutos estratégicos: depende de cómo se compita, como en la política comercial estratégica \vertema{3.B.8}.
\item \emph{El efecto disuasorio}: el exportador sube sus precios antes de la denuncia, para no provocarla.
\end{itemize}

\paragraph{d) Evidencia empírica}

\begin{itemize}
\item En Estados Unidos, \textcite{Prusa2001Spread} estima que un derecho antidumping reduce las importaciones de los países investigados entre un 30~y un 50~\%, y que las de los no investigados aumentan: buena parte de la protección se pierde por \emph{desviación del comercio}.
\item En la UE, \textcite{KoningsVandenbussche2001Diversion} encuentran que esa desviación es mucho menor que en Estados Unidos; una explicación es que la UE suele investigar a la vez a varios países y extiende las medidas a quien las elude.
\end{itemize}

\paragraph{e) Valoración}

\begin{itemize}
\item La defensa comercial es una válvula que sostiene la apertura, pero su coste lo pagan los consumidores y las industrias que usan esos insumos.
\item Su lógica es la del precio, y no responde a un socio que usa el comercio para presionar políticamente, que es lo que han traído los nuevos instrumentos.
\end{itemize}

% ==========================================================================
\section{Anexo: los acuerdos comerciales de la UE por regiones}\label{anexo:acuerdos}

\begin{table}[H]
\centering
\caption{Los acuerdos comerciales preferenciales de la UE en vigor o en aplicación provisional, noviembre de 2026}\label{tab:acuerdos}
\footnotesize
\begin{tabular}{>{\raggedright\arraybackslash}m{2.4cm}>{\raggedright\arraybackslash}m{5.6cm}>{\raggedright\arraybackslash}m{6.6cm}}
\toprule
\rowcolor{tceeazulpalido} \textbf{Región} & \textbf{Socios y tipo de acuerdo} & \textbf{Notas}\\
\midrule
Europa occidental & Espacio Económico Europeo (Noruega, Islandia, Liechtenstein, 1994); Suiza (ALC de 1972 y unos 120 acuerdos bilaterales) & El EEE da acceso al mercado interior a cambio de adoptar el acervo; el Tribunal de la AELC resuelve las diferencias. Suiza rechazó el EEE en referéndum (1992); el paquete firmado en 2026 ordena su relación\\
Uniones aduaneras & Turquía (1995, productos industriales), Andorra (1991), San Marino (2002), Mónaco & Turquía aplica el arancel común, pero no la política comercial: los socios de la UE no están obligados a abrirle su mercado\\
Balcanes y vecindad oriental & Acuerdos de estabilización y asociación (Albania, Bosnia y Herzegovina, Kosovo, Macedonia del Norte, Montenegro, Serbia); DCFTA con Ucrania, Moldavia y Georgia (2014-2016) & 9 países candidatos a la adhesión \vertema{3.B.36}\\
Reino Unido & Acuerdo de Comercio y Cooperación (2021) & Sin aranceles ni contingentes, con reglas de origen (anexo~\ref{anexo:reino-unido})\\
Mediterráneo & Acuerdos euromediterráneos de asociación (Marruecos, Argelia, Túnez, Egipto, Jordania, Líbano, Israel, Autoridad Palestina) & Libre comercio industrial y concesiones agrícolas recíprocas; con Siria, suspendido \vertema{3.B.42}\\
América del Norte & Canadá (CETA, 2017, aplicación provisional); México (2000; modernización firmada en 2026) & CETA: 18 de 27 ratificaciones en noviembre de 2026\\
América Latina y Caribe & Chile (2003; acuerdo interino desde 2025); Colombia, Perú (2013) y Ecuador (2017); Centroamérica (2013); Cariforum (2008); Mercosur (2026, provisional) & \\
Asia & Corea del Sur (2011); Japón (2019); Singapur (2019); Vietnam (2020) & Acuerdos de protección de inversiones con Singapur y Vietnam pendientes de ratificación\\
Oceanía & Nueva Zelanda (2024); Papúa Nueva Guinea (2014), Fiyi (2014), Samoa (2018) e Islas Salomón (2020) & \\
África & Acuerdos de asociación económica con la SADC (2016), África Oriental y Austral (2012), Camerún (2014), Costa de Marfil y Ghana (2016) y Kenia (2024) & Los acuerdos regionales de África Occidental y de la Comunidad de África Oriental, pendientes \vertema{3.B.42}\\
\bottomrule
\end{tabular}
\fuente{elaboración propia a partir de la Comisión Europea\nocite{CE2026Negociaciones,CE2026Informe}. Año de aplicación (provisional o definitiva).}
\end{table}

\begin{anotaciones}
\textbf{Para saber más} (enlaces comprobados en noviembre de 2026):
\begin{itemize}
\item Consejo de la UE: \href{https://www.consilium.europa.eu/es/policies/trade-policy/trade-agreements/}{los acuerdos comerciales de la UE}, y las infografías \href{https://www.consilium.europa.eu/es/infographics/the-eu-s-role-in-global-trade/}{el papel de la UE en el comercio mundial}, \href{https://www.consilium.europa.eu/es/infographics/trade-defence-instruments/}{los instrumentos de defensa comercial} y \href{https://www.consilium.europa.eu/en/infographics/eu-trade-map/}{el mapa de los acuerdos}.
\item Comisión Europea: \href{https://ec.europa.eu/trade/policy/eu-position-in-world-trade/}{la UE en el comercio mundial} y \href{https://ec.europa.eu/commission/presscorner/detail/es/ip_20_1867}{el control de las inversiones extranjeras (2020)}; Unión Europea: \href{https://european-union.europa.eu/principles-countries-history/joining-eu_es}{la adhesión y los países candidatos}.
\item OMC: \href{https://www.wto.org/spanish/tratop_s/adp_s/adp_info_s.htm}{información técnica sobre el antidumping}.
\item Prensa y divulgación: \href{https://www.eldiario.es/politica/eurocamara-aprueba-mayores-poderes-ue-controversias-comerciales_1_6944760.html}{la reforma del Reglamento de cumplimiento (eldiario.es, 2021)}; \href{https://www.iisd.org/itn/en/2020/10/05/ceta-faces-hurdle-after-cypriot-parliament-fails-to-ratify-the-agreement/}{el rechazo chipriota al CETA (IISD, 2020)}; \href{https://www.statista.com/chart/5100/uk-chooses-brexit/}{el resultado del referéndum del Brexit (Statista)}; \href{https://en.wikipedia.org/wiki/European_Free_Trade_Association}{la AELC y el EEE (Wikipedia)}; \href{https://www.youtube.com/watch?v=bVDjaTO_y9Y}{por qué Suiza no está en la UE (vídeo, 2022)}.
\item Manual: \textcite[cap.~6]{Camarero2019Economia}.
\end{itemize}
\end{anotaciones}

% ==========================================================================
\section{Anexo: el Acuerdo de Comercio\cortetitulo{} y Cooperación con el Reino Unido}\label{anexo:reino-unido}

\begin{itemize}
\item \emph{Relevancia}: antes del Brexit, el Reino Unido hacía el 13~\% del comercio de mercancías de la UE con terceros; en 2025 sigue siendo su tercer socio (9,7~\% del comercio de bienes y el segundo en servicios)\nocite{Eurostat2026Socios}.
\item \emph{Mercancías}:
  \begin{itemize}
  \item sin aranceles ni contingentes para los productos que cumplen las reglas de origen (un automóvil británico con batería china puede no cumplirlas); sin el acuerdo, se habrían aplicado los aranceles NMF: un 10~\% a los automóviles, un 12~\% a la ropa y al calzado y más del 40~\% a muchos productos agroalimentarios;
  \item hay controles aduaneros y sanitarios, porque no hay reconocimiento mutuo de las normas; los operadores económicos autorizados tienen facilidades;
  \item anexos sectoriales para el vino, los automóviles, la química y los medicamentos;
  \item en Irlanda del Norte se siguen aplicando las normas de mercancías de la UE (Protocolo, ajustado por el Marco de Windsor en 2023), para evitar una frontera en la isla.
  \end{itemize}
\item \emph{Servicios}: el Reino Unido pierde la libre prestación y el pasaporte financiero; hay una cláusula NMF hacia el futuro (salvo en los servicios financieros), facilidades para los viajes de negocios (modo 4) y para los abogados, pero no reconocimiento automático de cualificaciones. Las equivalencias financieras son decisiones unilaterales de cada parte.
\item \emph{Transporte}: las aerolíneas británicas tienen las libertades tercera y cuarta, pero no pueden hacer rutas dentro de la UE; las europeas tienen que seguir siendo de propiedad y control europeos, lo que obligó a cambios en el accionariado de algunas.
\item \emph{Otros}: protección de las indicaciones geográficas registradas hasta 2020, contratación pública por encima del Acuerdo de la OMC, comercio digital, energía (el Reino Unido sale del mercado de derechos de emisión de la UE), pesca, participación en Horizonte Europa y cooperación policial y judicial. Salió de Erasmus+, al que vuelve en 2027.
\item \emph{Gobernanza}: un Consejo de Asociación, un mecanismo de solución de diferencias por arbitraje y la posibilidad de \emph{reequilibrar} con aranceles si las normas de una parte divergen mucho de las de la otra.
\end{itemize}

\clearpage
% ==========================================================================
\section{Anexo: la negociación con Estados Unidos\cortetitulo{} con la solución de \autor{Nash}}\label{anexo:nash}

\begin{anotaciones}
\textbf{Notación.} $u_i$ es la ganancia de bienestar de cada parte, $d_i$ su punto de desacuerdo (lo que obtiene sin acuerdo) y $F$ la frontera de los acuerdos posibles.
\end{anotaciones}

\paragraph{a) Idea}

\begin{itemize}
\item Un país grande puede ganar con su arancel óptimo, pero si el otro responde, la guerra arancelaria deja a los 2~peor que el libre comercio: es el equilibrio de \autor{Nash} en aranceles de \textcite{Johnson1953Optimum} \vertema{3.B.7} \vertema{3.B.8}. Las reglas de la OMC están para evitarla \vertema{3.B.33}.
\item Cuando las reglas no se respetan, los países negocian, y el reparto lo decide lo que cada uno pierde si no hay acuerdo. Es la lógica de la solución de negociación de \textcite{Nash1950Bargaining}, que \textcite{GrossmanHelpman1995TradeWars} aplican a las guerras y las negociaciones comerciales.
\end{itemize}

\paragraph{b) Modelo}

\subparagraph{i) Supuestos}

\begin{itemize}
\item 2~partes, la UE y Estados Unidos, negocian un acuerdo; $u_{\mathit{UE}}$ y $u_{\mathit{EEUU}}$ son sus ganancias de bienestar.
\item $F$ es el conjunto de acuerdos posibles, convexo, con una frontera decreciente: lo que gana uno de más, lo pierde el otro.
\item El \concepto{punto de desacuerdo} $d=(d_{\mathit{UE}},d_{\mathit{EEUU}})$ es lo que obtiene cada parte si no hay acuerdo: la guerra arancelaria con represalias, que deja a cada uno con una pérdida que depende de cuánto exporta al otro y de qué más está en juego (la seguridad, por ejemplo).
\item La solución cumple los axiomas de \autor{Nash}: eficiencia, simetría, invariancia a las transformaciones lineales de la utilidad e independencia de las alternativas irrelevantes.
\end{itemize}

\subparagraph{ii) Desarrollo}

\begin{itemize}
\item \emph{Paso 1: la solución de Nash}.
  \begin{itemize}
  \item El acuerdo maximiza el producto de las ganancias de cada parte sobre su punto de desacuerdo:
  \funcionclave[eq:nash]{\optim{\max}{u_{\mathit{UE}},\,u_{\mathit{EEUU}}}{\big(u_{\mathit{UE}}-d_{\mathit{UE}}\big)\big(u_{\mathit{EEUU}}-d_{\mathit{EEUU}}\big)}{(u_{\mathit{UE}},u_{\mathit{EEUU}})\in F}}
  \item En el óptimo, la pendiente de la frontera es igual a la de la recta que une $d$ con el acuerdo: $-\dfrac{du_{\mathit{EEUU}}}{du_{\mathit{UE}}}=\dfrac{u_{\mathit{EEUU}}-d_{\mathit{EEUU}}}{u_{\mathit{UE}}-d_{\mathit{UE}}}$.
  \end{itemize}
\item \emph{Paso 2: un desacuerdo simétrico}.
  \begin{itemize}
  \item Con la frontera $u_{\mathit{EEUU}}=10-u_{\mathit{UE}}^2/10$ y $d=(0,0)$, el acuerdo es $N=(5{,}8;\ 6{,}7)$ (gráfico~\ref{fig:negociacion}).
  \end{itemize}
\item \emph{Paso 3: la UE pierde más sin acuerdo}.
  \begin{itemize}
  \item Si la ruptura le cuesta más a la UE (exporta a Estados Unidos mucho más de lo que le compra en bienes, y depende de él en defensa), su punto de desacuerdo baja a $d'=(-4,0)$, y el acuerdo se desplaza a $N'=(4{,}6;\ 7{,}9)$: menos para la UE.
  \end{itemize}
\item \emph{Paso 4: las represalias creíbles}.
  \begin{itemize}
  \item Si la UE puede hacer daño a Estados Unidos (contramedidas preparadas, el Instrumento Anticoerción, los servicios digitales), el desacuerdo de Estados Unidos también baja, a $d''=(-4,-4)$, y el acuerdo vuelve hacia la UE: $N''=(5{,}6;\ 6{,}8)$.
  \end{itemize}
\end{itemize}

\begin{figure}[H]
\centering
\caption{La negociación entre la UE y Estados Unidos con la solución de \autor{Nash}}\label{fig:negociacion}
\grafico[0.7\textwidth]{negociacion}
\fuente{elaboración propia a partir de \textcite{Nash1950Bargaining} y \textcite{GrossmanHelpman1995TradeWars}, con el ejemplo del texto ($u_{\mathit{EEUU}}=10-u_{\mathit{UE}}^2/10$).}
\end{figure}

\subparagraph{iii) Implicaciones}

\begin{itemize}
\item \emph{Gana más quien menos pierde con la ruptura}: el reparto mejora para la UE si sube su punto de desacuerdo y empeora si sube el del otro.
\item \emph{Explica el acuerdo asimétrico de 2025}: la UE aceptó un 15~\% frente a un arancel cero propio porque, en una guerra comercial, perdía más que Estados Unidos y porque la negociación estaba unida a la seguridad de Europa.
\item \emph{Y el valor de la amenaza}: las contramedidas y el Instrumento Anticoerción no sirven tanto para usarse como para mejorar el punto de desacuerdo. Es «donde las dan, las toman» como disuasión: la sola voz de \autor{Meunier} (\ref{sec:meunier}) solo pesa si detrás hay capacidad de responder.
\end{itemize}

\paragraph{c) Extensiones}

\begin{itemize}
\item \textcite{GrossmanHelpman1995TradeWars} añaden los grupos de presión: cada gobierno pesa más a los productores organizados, y por eso los acuerdos protegen los sectores con más influencia (el automóvil en la UE, el acero en Estados Unidos).
\item En un juego repetido, la cooperación se sostiene con la amenaza de volver a la guerra; las reglas de la OMC, con su reciprocidad y su no discriminación, son la forma de hacerla creíble y de acotarla \vertema{3.B.33}.
\end{itemize}

\paragraph{d) Evidencia empírica}

\begin{itemize}
\item Los aranceles de 2018-2019 los pagaron casi enteros los compradores estadounidenses: el precio sin arancel de las importaciones apenas bajó\nocite{Amiti2019Aranceles,Fajgelbaum2020Proteccionismo}. Es la señal de que Estados Unidos no mejoró su RRI y de que la guerra tiene costes para quien la empieza, lo que sube su punto de desacuerdo menos de lo que parece.
\item Los resultados de 2025 encajan con el modelo: el Reino Unido, con un comercio de bienes más equilibrado con Estados Unidos, obtuvo un 10~\%; la UE y Japón, con grandes superávits, un 15~\%.
\end{itemize}

\paragraph{e) Valoración}

\begin{itemize}
\item El modelo explica el reparto, pero no que el acuerdo sea estable: es una declaración política, no un tratado ni un acuerdo del art. XXIV, y Estados Unidos ha cambiado 2~veces su base legal en 2026.
\item Y es estático: a largo plazo, el mejor punto de desacuerdo de la UE es tener alternativas, que es lo que construye su red de acuerdos con el resto del mundo.
\end{itemize}

% ==========================================================================
\section{Anexo: preguntas de test}\label{anexo:test}

\begin{anotaciones}
Preguntas tipo test del primer ejercicio: \href{https://www.victorgutierrezmarcos.es/oposicion/temario/primer-ejercicio/test/simulador.html?tema=3.B.41}{simulador de test}.
\end{anotaciones}

\bibliografia
\end{document}
